%%%%%%%%%%%%%%%%%%%%%%%%%%%%%%%%%%%%%%%%%%%%%%%%%%%%%%%%%%%%%%%%%%%%%%
% Lớp 10 | Mệnh đề. Tập hợp | Mệnh đề
% Dạng 0D1?1-5: Mệnh đề với mọi, tồn tại (và phủ định chúng)
% Số câu: 101  (sinh tự động bởi tach_cau.py; từ nay thầy sửa trực tiếp file này)
%%%%%%%%%%%%%%%%%%%%%%%%%%%%%%%%%%%%%%%%%%%%%%%%%%%%%%%%%%%%%%%%%%%%%%

%%%%%%%%%%%%%%%% TRẮC NGHIỆM 4 PHƯƠNG ÁN %%%%%%%%%%%%%%%%

% nguồn: MaHoaNH_TracNghiem_Lop10
\begin{ex}%[0D1N1-5]%[ma:608e3343]
	Mệnh  đề \lq\lq $\forall x\in \mathbb{R},\,\,x^2-x+3<0$ \rq\rq\, có mệnh đề phủ định là
	\choice 
		{\True \lq\lq $\exists x  \in \mathbb{R},\,\,x^2-x+3\ge 0$\rq\rq}
		 {\lq\lq $\forall x \in \mathbb{R},\,\,x^2-x+3>0$\rq\rq}
		 {\lq\lq $\forall x \in \mathbb{R},\,\,x^2-x+3<0$\rq\rq}
		{\lq\lq $\nexists  x \in \mathbb{R},\,\,x^2-x+3<0$\rq\rq}
	\loigiai{ 
	\lq\lq $\exists  x \in \mathbb{R},\,\,x^2-x+3\ge 0$ \rq\rq.}
\end{ex}

% nguồn: MaHoaNH_TracNghiem_Lop10
\begin{ex}%[0D1N1-5]%[ma:811f9ff5]
	Trong các mệnh đề sau, mệnh đề nào đúng? 
	 \choice
		{$\forall n\in  \mathbb{N}, n > 1$}
		{\True $\exists n\in  \mathbb{Q}, n^2=n$}
		{$\exists n\in  \mathbb{N}, n\left(n+1\right)$ là một số lẻ}
		{$\forall n\in \mathbb{N}, n^2>  n$}
	\loigiai{ 
		\begin{itemize}
			\item Với $n=0$, thì mệnh đề  \lq\lq $\forall n\in \mathbb{N}, n > 1$\rq\rq sai.
			\item Với $n=0$, thì mệnh đề \lq\lq$\exists n\in \mathbb{Q}, n^2=n$ \rq\rq  đúng.
			 \item Với $n\in \mathbb{N}$ thì $n\left(n+1\right)$ là một số chẵn.
			\item Với $n=0$, thì mệnh đề \lq\lq $\forall n\in \mathbb{N}, n^2>  n$ \rq\rq sai.
		\end{itemize}
	} 
\end{ex}

% nguồn: MaHoaNH_TracNghiem_Lop10
\begin{ex}%[0D1H1-5]%[ma:9e3b3327]
	Mệnh đề phủ định của mệnh đề “$\forall  x\in\mathbb{R},x^2-x+2>0$” là mệnh đề nào dưới đây?
	 \choice
		{$\exists  x\in\mathbb{R},\,x^2-x+2<0$}
		{\True $\exists x\in\mathbb{R},\,x^2-x+2\le  0$}
		{$\forall  x\in\mathbb{R},\,x^2-x+2\le 0$}
		{$\forall  x\in\mathbb{R},\,x^2-x+2<0$}
	 \loigiai{
	Phủ định của mệnh đề $\forall x\in X,\, P(x)$ là  mệnh đề $\exists x\in X,\,\overline{P(x)}$.}
\end{ex}

% nguồn: MaHoaNH_TracNghiem_Lop10
\begin{ex}%[0D1H1-5]%[ma:2ea0ced3]
	Cho mệnh đề \lq\lq Có một học sinh trong lớp C4 không chấp hành luật  giao thông\rq\rq. Mệnh đề phủ định của mệnh đề này là
	\choice 
		{Không có học sinh nào trong lớp C4 chấp hành luật  giao thông}
		{\True Mọi học sinh  trong lớp C4 đều chấp hành luật giao thông}
		{Có một học sinh  trong lớp C4 chấp hành luật giao thông}
		{Mọi học sinh trong lớp C4 không  chấp hành luật giao thông}
	\loigiai{ 
		Mệnh đề phủ định  là \lq\lq Mọi học sinh trong lớp C4 đều chấp hành luật giao thông\rq\rq.
	} 
\end{ex}

% nguồn: MaHoaNH_TracNghiem_Lop10
\begin{ex}%[0D1H1-5]%[ma:135d3b70]
	Mệnh đề nào sau  đây đúng?
	\choice 
		{\True $\forall x \in \mathbb{R}, x^2-x+1>0$} 
		{$\exists n \in  \mathbb{N}, n<0$}
		{$\exists x  \in \mathbb{Q}, x^2=2$}
		 {$\forall x \in \mathbb{Z}, \dfrac{1}{x}>0$}
	\loigiai{ 
		Vi $x^2-x+1=\left(x-\dfrac{1}{2}\right)^2+\dfrac{3}{4}>0, \forall x  \in \mathbb{R}$.
	 }
\end{ex}

% nguồn: MaHoaNH_TracNghiem_Lop10
\begin{ex}%[0D1H1-5]%[ma:06f4dc66]
	 Mệnh đề \lq\lq Mọi động vật đều đi chuyển\rq\rq \,có mệnh đề phủ định là
	\choice 
		{Có ít nhất một động vật di chuyển} 
		 {Mọi động vật đều đứng yên}
		{\True Có  ít nhất một động vật không di chuyển}
		{Mọi  động vật đều không di chuyển}
	\loigiai{ 
		Mệnh đề phủ định là \lq\lq Có ít nhất một động  vật không di chuyển\rq\rq.
	} 
\end{ex}

% nguồn: MaHoaNH_TracNghiem_Lop10
\begin{ex}%[0D1H1-5]%[ma:a6960496]
	Tìm  mệnh đề đúng.
	\choice 
		{\lq\lq  $\exists x; x^2+3=0$ \rq\rq}
		{\lq\lq $\exists x;  x^4+3x^2+2=0 $\rq\rq}
		{\True \lq\lq $\forall x \in \mathbb{Z}; x^5>\mathrm{x}^2$  \rq\rq}
		{\lq\lq $\forall n \in  \mathbb{N};\left((2 n+1)^2-1\right) \vdots 4$ \rq\rq}
	\loigiai{ 
		Ta biến đổi $(2 n+1)^2-1=4 n^2+4 n=4\left(n^2+n\right)  \vdots 4; \forall n \in \mathbb{N}$.
	 }
\end{ex}

% nguồn: MaHoaNH_TracNghiem_Lop10
\begin{ex}%[0D1H1-5]%[ma:11c14128]
	 Chọn mệnh đề sai.
	 \choice
		{\True \lq\lq $\forall x \in \mathbb{R}\colon x^2>0 $\rq\rq} 
		{\lq\lq $\exists n \in  \mathbb{N}: n^2=n$ \rq\rq}
		 {\lq\lq $\forall n \in \mathbb{N}: n \leq 2 n$ \rq\rq}
		 {\lq\lq $\exists x \in \mathbb{R}\colon x<1$ \rq\rq}
	 \loigiai{
		Với $x=0 \in \mathbb{R}$ thì $x^2=0$  nên \lq\lq $\forall x \in \mathbb{R}\colon x^2>0$ \rq\rq\, sai.
	} 
\end{ex}

% nguồn: MaHoaNH_TracNghiem_Lop10
\begin{ex}%[0D1H1-5]%[ma:b103d026]
	Tìm mệnh đề phủ định của mệnh đề \lq\lq  $\forall x \in \mathbb{R}\colon x^2>x$ \rq\rq.
	 \choice
		{$\forall  x \in \mathbb{R}\colon x^2 \leq x$}
		{$\exists x \in \mathbb{R}\colon  x^2>x$}
		{\True  $\exists x \in \mathbb{R}\colon x^2 \leq x$}
		{$\exists  x \in \mathbb{R}\colon x^2<x$}
	\loigiai{ 
		Mệnh đề $A$ : \lq\lq $\forall x  \in \mathbb{R}\colon x^2>x$ \rq\rq\, phủ định là \lq\lq $\exists x \in \mathbb{R}\colon x^2 \leq x$ \rq\rq.
	} 
\end{ex}

% nguồn: MaHoaNH_TracNghiem_Lop10
\begin{ex}%[0D1H1-5]%[ma:10d8dac9]
	Cho mệnh đề: \lq\lq  $\forall x \in \mathbb{R}, x^2+3x+5>0$ \rq\rq. Mệnh đề phủ định của mệnh đề trên là
	 \choice
		{$\forall x \in \mathbb{R}, x^2+3x+5 \leq  0$}
		{\True $\exists x  \in \mathbb{R}, x^2+3x+5 \leq 0$}
		{$\forall x \in \mathbb{R}, x^2+3x+5<0$} 
		{$\exists x  \in \mathbb{R}, x^2+3x+5>0$}
	\loigiai{ 
		Chú ý: Phủ định của mệnh đề \lq\lq $\forall x \in X, P(x)$ \rq\rq là \lq\lq $\exists x \in X, \overline{P(x)}$ \rq\rq . 
	 }
\end{ex}

% nguồn: MaHoaNH_TracNghiem_Lop10
\begin{ex}%[0D1H1-5]%[ma:049d4ab0]
	Mệnh đề  phủ định của mệnh đề: $\forall x \in \mathbb{R}, x^2+x+5 > 0$ là
	\choice 
		{\True $\exists x  \in \mathbb{R}, x^2+x+5 \leq 0$}
		{$\forall x \in \mathbb{R}, x^2+x+5 \leq  0$}
		{$\exists x \in  \mathbb{R}, x^2+x+5<0$}
		 {$\forall x \in \mathbb{R}, x^2+x+5<0$}
	 \loigiai{
		Mệnh đề phủ định của mệnh  đề: $\forall x \in \mathbb{R}, x^2+x+5 > 0$ là $\exists x \in \mathbb{R}, x^2+x+5 \leq 0$.
	 }
\end{ex}

% nguồn: MaHoaNH_TracNghiem_Lop10
\begin{ex}%[0D1H1-5]%[ma:1638071e]
	Chọn mệnh đề đúng. 
	\choice 
		{$\forall n\in {\bf \mathbb{N}}^{*}, n^2-1$ là bội số  của $3$}
		{$\exists x\in \mathbb{Q},  x^2=3$}
		{$\forall n\in {\bf \mathbb{N}}, 2^n+1$ là số nguyên tố} 
		{\True $\exists n\in {\bf  \mathbb{N}}, 2^n \ge n+2$}
	 \loigiai{
		Mệnh đề $\exists n\in {\bf \mathbb{N}}, 2^n \ge  n+2$ đúng vì tồn tại $2\in {\bf \mathbb{N}}, 2^2 \ge 2+2$.
	 }
\end{ex}

% nguồn: MaHoaNH_TracNghiem_Lop10
\begin{ex}%[0D1H1-5]%[ma:67b9826e]
	Mệnh đề phủ định  của mệnh đề \lq\lq $\exists x \in \mathbb{R}, x^2+x+13=0$ \rq\rq\, là
	\choice 
		{\True \lq\lq $\forall x \in  \mathbb{R}, x^2+x+13 \neq 0$ \rq\rq}
		{\lq\lq  $\exists x \in \mathbb{R}, x^2+x+13>0$ \rq\rq}
		 {\lq\lq $\forall x \in \mathbb{R}, x^2+x+13=0$ \rq\rq}
		{\lq\lq $\exists x \in \mathbb{R},  x^2+x+13 \neq 0$ \rq\rq}
	\loigiai{ 
		Mệnh đề  phủ định của mệnh đề \lq\lq $\exists x \in \mathbb{R}, x^2+x+13=0$ \rq\rq\, là \lq\lq $\forall x \in \mathbb{R}, x^2+x+13 \neq 0$ \rq\rq.
	} 
\end{ex}

% nguồn: MaHoaNH_TracNghiem_Lop10
\begin{ex}%[0D1H1-5]%[ma:8a5e85b2]
	Cho mệnh đề: \lq\lq Có một học sinh trong lớp 10A không thích học môn Toán\rq\rq. Mệnh đề phủ  định của mệnh đề này là
	 \choice
		{\True \lq\lq Mọi học sinh  trong lớp 10A đều thích học môn Toán\rq\rq}
		{\lq\lq Mọi  học sinh trong lớp 10A đều không thích học môn Toán\rq\rq}
		{\lq\lq Mọi học sinh  trong lớp 10A đều thích học môn Văn\rq\rq}
		{\lq\lq Có một học sinh trong lớp 10A thích học  môn Toán\rq\rq}
	\loigiai{ 
		Mệnh đề  phủ định là \rq\rq Mọi học sinh trong lớp 10A đều thích học môn Toán\rq\rq.
	 }
\end{ex}

% nguồn: MaHoaNH_TracNghiem_Lop10
\begin{ex}%[0D1H1-5]%[ma:66fa32a1]
	Phủ định của mệnh đề: \lq\lq $\forall n \in \mathbb{N}:  n^2+1$ không chia hết cho $3$ \rq\rq là
	\choice 
		 {\lq\lq $\forall n \in \mathbb{N}: n^2+1$ chia hết cho $3$ \rq\rq}
		{\lq\lq  $\exists n \in \mathbb{N}: n^2+1$ không chia hết cho $3$ \rq\rq}
		{\True \lq\lq $\exists n \in \mathbb{N}: n^2+1$ chia  hết cho $3$ \rq\rq}
		{\lq\lq $\nexists n \in \mathbb{N}:  n^2+1$ không chia hết cho $3$ \rq\rq}
	 \loigiai{
		Phủ định của $\forall$ là $\exists$.\\ 
		Phủ định của \lq\lq không  chia hết\rq\rq\,là \,\lq\lq chia hết.\rq\rq
	 }
\end{ex}

% nguồn: MaHoaNH_TracNghiem_Lop10
\begin{ex}%[0D1H1-5]%[ma:e31956d1]
	Phủ định của mệnh đề: \lq\lq $\exists x \in \mathbb{R}\colon x^2+x+1$ là số dương\rq\rq  \,là
	 \choice
		 {\True \lq\lq $\forall x \in \mathbb{R}\colon x^2+x+1$ là số không dương\rq\rq}
		{\lq\lq  $\forall x \in \mathbb{R}\colon x^2+x+1$ là số âm\rq\rq}
		{\lq\lq  $\forall x \in \mathbb{R}\colon x^2+x+1$ là số dương\rq\rq}
		{\lq\lq $\nexists x \in \mathbb{R}\colon x^2+x+1$ là số dương\rq\rq} 
	\loigiai{ 
		Phủ định của $\exists$ là $\forall$.\\ 
		Phủ định  của \lq\lq số dương\rq\rq\,là \lq\lq số không dương\rq\rq.
	} 
\end{ex}

% nguồn: MaHoaNH_TracNghiem_Lop10
\begin{ex}%[0D1H1-5]%[ma:a347dc20]
	Cho mệnh đề  \lq\lq $\forall x \in \mathbb{R}, x^2-x+7<0$ \rq\rq. Hỏi mệnh đề nào là mệnh đề phủ định của mệnh đề trên?
	 \choice
		{\True $\exists x \in \mathbb{R}, x^2-x+7 \geq  0$}
		{$\forall x \in \mathbb{R}, x^2-x+7>0$} 
		{$\forall x \in \mathbb{R}, x^2-x+7<0$} 
		{$\exists x  \in \mathbb{R}, x^2-x+7<0$}
	 \loigiai{
		 $\forall x \in \mathbb{R}, x^2-x+7>0$ sai là vì không dùng đúng kí hiệu của phủ định.\\
		$\forall x \in \mathbb{R}, x^2-x+7<0$ sai  là vì không dùng đúng $\geq$.\\
		$\exists x \in  \mathbb{R}, x^2-x+7<0$ chưa phủ định vế sau của mệnh đề.
	 }
\end{ex}

% nguồn: MaHoaNH_TracNghiem_Lop10
\begin{ex}%[0D1H1-5]%[ma:d56498c5]
	Mệnh đề phủ định  của mệnh đề \lq\lq Phương trình $a x^2+b x+c=0\,(a \neq 0)$ vô nghiệm\rq\rq\, là mệnh đề nào sau đây?
	\choice 
		{\True  Phương trình $a x^2+b x+c=0\,(a \neq 0)$ có nghiệm}
		{Phương trình $a x^2+b x+c=0\,(a  \neq 0)$ có $2$ nghiệm phân biệt}
		 {Phương trình $a x^2+b x+c=0\,(a \neq 0)$ có nghiệm kép}
		{Phương trình $a x^2+b x+c=0\,(a  \neq 0)$ không có nghiệm}
	 \loigiai{
		Phương án phương trình $a x^2+b x+c=0\,(a \neq 0)$ có nghiệm đúng vì phủ định vô nghiệm là có  nghiệm.\\
		Phương án phương trình $a x^2+b x+c=0\,(a \neq 0)$ có $2$ nghiệm phân biệt sai vì học sinh nhầm phủ định vô nghiệm là phương trình sẽ có $2$ nghiệm phân biệt.\\ 
		Đáp án Phương trình $a x^2+b x+c=0\,(a \neq 0)$ có nghiệm kép sai vì học sinh nhầm phủ định vô nghiệm là có $1$ nghiệm tức nghiệm  kép.\\
		Phương án phương trình $a x^2+b x+c=0\,(a \neq 0)$ không có nghiệm sai vì học sinh không hiểu câu hỏi của đề, học sinh  nghĩ vô nghiệm là không có nghiệm.
	} 
\end{ex}

% nguồn: MaHoaNH_TracNghiem_Lop10
\begin{ex}%[0D1H1-5]%[ma:5644c216]
	Cho mệnh  đề: \lq\lq $\exists x \in \mathbb{R} \mid 2x^2-3x-5<0$ \rq\rq. Mệnh đề phủ định sẽ là
	 \choice
		{\True \lq\lq $\forall x \in \mathbb{R}, 2x^2+3x-5  \geq 0$ \rq\rq}
		{\lq\lq $\forall  x \in \mathbb{R}, 2x^2+3x-5>0$ \rq\rq}
		{\lq\lq  $\exists x \in \mathbb{R}, 2x^2+3x-5>0$ \rq\rq}
		{\lq\lq $\exists x \in \mathbb{R}, 2x^2+3x-5 \geq 0$  \rq\rq}
	\loigiai{ 
		Phương án \lq\lq $\forall x \in \mathbb{R}, 2x^2+3x-5 \geq 0$ \rq\rq\, đúng vì phủ định của \lq\lq $\exists$  \rq\rq \,là \lq\lq $\forall$ \rq\rq \,và phủ định của dấu \lq\lq $<$ \rq\rq \,\,là dấu \lq\lq $\geq$ \rq\rq.\\
		Phương án \lq\lq $\forall x \in \mathbb{R}, 2x^2+3x-5>0$ \rq\rq sai vì học sinh nhầm phủ định của  dấu \lq\lq <\rq\rq \,là dấu \lq\lq >\rq\rq.\\
		Phương án \lq\lq $\exists x \in \mathbb{R}, 2x^2+3x-5>0$ \rq\rq \,sai vì học sinh không nhớ phủ định của \lq\lq $\exists$ \rq\rq \,là \lq\lq $\forall$ \rq\rq \,và phủ định dấu \lq\lq <\rq\rq\, là dấu \lq\lq  $\geq$ \rq\rq.\\
		Phương án \lq\lq $\exists x \in \mathbb{R}, 2x^2+3x-5 \geq 0$ \rq\rq\, sai vì học sinh không nhớ phủ định của \lq\lq $\exists$ \rq\rq \,là \lq\lq $\forall$  \rq\rq.
	 }
\end{ex}

% nguồn: MaHoaNH_TracNghiem_Lop10
\begin{ex}%[0D1H1-5]%[ma:3d3cdf08]
	Trong các mệnh đề sau, mệnh đề nào sai? 
	\choice 
		{ $\exists n \in \mathbb{N}, n^2+11 n+2$  chia hết cho $11$}
		{\True $\exists  n \in \mathbb{N}, n^2+1$ chia hết cho $4$}
		{Tồn  tại số nguyên tố chia hết cho $5$}
		{$\exists  n \in \mathbb{Z}, 2x^2-8=0$}
	\loigiai{ 
		Khi $n=3$ thì giá trị của $\left(n^2+11 n+2\right)$ bằng  $44: 11$ nên phương án $\exists n \in \mathbb{N}, n^2+11 n+2$ chia hết cho $11$ đúng.\\
		Khi $n=2 k, k \in  \mathrm{~N} \Rightarrow n^2+1=4 k^2+1$ không chia hết cho $4, k \in \mathrm{~N}$.\\
		Khi $n=2 k+1, k \in  \mathrm{~N} \Rightarrow n^2+1=(2 k+1)^2+1=4 k^2+4 k+2$ không chia hết cho $4, k \in \mathrm{~N}$.\\
		Phương án \lq\lq Tồn tại số nguyên tố chia  hết cho $5$\rq\rq\, đúng.\\
		Phương trình $2x^2-8=0 \Leftrightarrow x^2=4 \Leftrightarrow x=-2; x=2 \in \mathbb{Z}$ nên phương án \lq\lq$\exists n \in  \mathbb{Z}, 2x^2-8=0$\rq\rq\, đúng.
	 }
\end{ex}

% nguồn: MaHoaNH_TracNghiem_Lop10
\begin{ex}%[0D1H1-5]%[ma:9a292bba]
	Trong các mệnh đề dưới đây, mệnh đề nào  là mệnh đề đúng?
	 \choice
		{\lq\lq $\exists x\in \mathbb{R}\mid x^2+1=0$\rq\rq} 
		{\lq\lq $\exists x\in  \mathbb{R}\mid x^2< 0$\rq\rq}
		{\True \lq\lq $\exists x\in \mathbb{N}\mid 2x^2-1<  0$\rq\rq}
		{\lq\lq $\exists x\in \mathbb{N}\mid x^2-2=0$\rq\rq} 
	 \loigiai{
		Ta  có $x^2\ge 0\Leftrightarrow x^2+1\ge 1$, $\forall x\in \mathbb{R}$.
		Nên  $\exists x\in \mathbb{R}\mid x^2+1=0$ là mệnh đề sai.\\
		Ta có $x^2\ge 0$, $\forall x\in \mathbb{R}$. Nên $\exists x\in \mathbb{R}\mid x^2< 0$ là  mệnh đề sai.\\
		$2x^2-1< 0\Leftrightarrow x^2< \dfrac{1}{2} \Leftrightarrow-\dfrac{\sqrt{2}}{2} < x < \dfrac{\sqrt{2}}{2}$, mà $x\in \mathbb{N}\Rightarrow x=0$. Vậy $\exists x\in \mathbb{N}\mid 2x^2-1< 0$ là mệnh  đề đúng.\\
		$x^2-2=0\Leftrightarrow x=\pm \sqrt{2} \notin \mathbb{N}$. Vậy $\exists x\in  \mathbb{N}\mid x^2-2=0$ là mệnh đề sai.
	 }
\end{ex}

% nguồn: MaHoaNH_TracNghiem_Lop10
\begin{ex}%[0D1K1-5]%[ma:56ec9c34]
	 Chọn mệnh đề đúng.
	\choice{\True $\forall x\in\mathbb{R}, x^2>5\Rightarrow  x>\sqrt{5}$ hoặc $x<-\sqrt{5}$}
		{$\forall x\in\mathbb{R},  x^2>5\Rightarrow -\sqrt{5}<x<\sqrt{5}$}
		 {$\forall x\in\mathbb{R}, x^2>5\Rightarrow x>\pm\sqrt{5}$}
		{$\forall x\in\mathbb{R}, x^2\ge5\Rightarrow x> \sqrt{5}$ hoặc $x<  -\sqrt{5}$}
	 \loigiai{Ta có $x^2 > 5 \Rightarrow x > \sqrt{5}$ hoặc $x < - \sqrt{5}$ với mọi $x \in \mathbb{R}$.}
\end{ex}

% nguồn: MaHoaNH_TracNghiem_Lop10
\begin{ex}%[0D1K1-5]%[ma:604257fe]
	Chọn mệnh đề đúng. 
	\choice{\True  $\forall x\in\mathbb{R}, x>3\Rightarrow x^2>9$}
		{$\forall x\in\mathbb{R}, x>-3\Rightarrow x^2>9$} 
		{$\forall  x\in\mathbb{R}, x^2>9\Rightarrow x>3$}
		{$\forall x\in\mathbb{R}, x^2>9\Rightarrow  x>-3$}
	\loigiai{Khi x $>3>0$  thì $x^2>3^2$ hay $x^2>9$.}
\end{ex}

% nguồn: MaHoaNH_TracNghiem_Lop10
\begin{ex}%[0D1K1-5]%[ma:e94495ea]
	Chọn mệnh  đề đúng.
	\choice{$\forall x\in\mathbb{R}, x^2\leq 16\Leftrightarrow  x\leq \pm4$}
		{\True $\forall x\in\mathbb{R}, x^2\leq  16\Leftrightarrow -4\leq x\leq 4$}
		{$\forall x\in\mathbb{R}, x^2\leq 16\Leftrightarrow  x\leq -4$ hoặc $x\geq 4$}
		{$\forall x\in\mathbb{R}, x^2\leq 16\Leftrightarrow  -4<x<4$}
	\loigiai{Ta có $x^2 \le 4 \Leftrightarrow -4 \le  x \le 4$ với $\forall x\in\mathbb{R}$.}
\end{ex}

% nguồn: MaHoaNH_TracNghiem_Lop10
\begin{ex}%[0D1B1-5]%[ma:c80aef44]
	Phủ định của mệnh đề ``$\exists x\in\mathbb{R}$, $5x-3x^2=1$'' là: 
	\choice 
		{``$\exists  x\in\mathbb{R}$, $5x-3x^2$''}
		{``$\forall x\in\mathbb{R}$, $5x-3x^2=1$''} 
		{``$\exists  x\in\mathbb{R}$, $5x-3x^2\geq 1$''}
		{\True ``$\forall x\in\mathbb{R}$, $5x-3x^2\ne  1$''}
	\loigiai{Phủ định của mệnh đề ``$\exists x\in\mathbb{R}$, $5x-3x^2=1$'' là ``$\forall x\in\mathbb{R}$,  $5x-3x^2\ne 1$.''}
\end{ex}

% nguồn: MaHoaNH_TracNghiem_Lop10
\begin{ex}%[0D1B1-5]%[ma:f5b4d6e2]
	Mệnh đề phủ định của mệnh đề $P$: “$\exists x\in\mathbb{R}:x-3>0$”  là:
	 \choice{$\overline{P}$: “$\exists x\in\mathbb{R}:x-3\leq 0$”}
		{\True $\overline{P}$:  “$\forall x\in\mathbb{R}:x-3\leq 0$”}
		{$\overline{P}$: “$\forall  x\in\mathbb{R}:x-3>0$”}
		{$\overline{P}$:  “$\exists x\notin\mathbb{R}:x-3>0$”}
	 \loigiai{
	Mệnh đề phủ định  là $\overline{P}$: “$\forall x\in\mathbb{R}:x-3\leq 0$”.}
\end{ex}

% nguồn: MaHoaNH_TracNghiem_Lop10
\begin{ex}%[0D1B1-5]%[ma:0f5e6653]
	Mệnh đề phủ  định của mệnh đề $P$: “$\forall x\in\mathbb{R}:x^2\geq 0$” là:
	\choice{$\overline{P}$: “$\exists x\in\mathbb{R}:x^2\leq 0$”} 
		 {$\overline{P}$: “$\forall x\in\mathbb{R}:x^2\leq 0$”}
		{\True  $\overline{P}$: “$\exists x\in\mathbb{R}:x^2<0$”}
		 {$\overline{P}$: “$\forall x\notin\mathbb{R}:x^2\geq 0$”}
	 \loigiai{
	Mệnh đề phủ định  là $\overline{P}$:“$\exists x\in\mathbb{R}:x^2<0$”.}
\end{ex}

% nguồn: MaHoaNH_TracNghiem_Lop10
\begin{ex}%[0D1B1-5]%[ma:ca8f23aa]
	Cho mệnh đề ``$\forall x \in  \mathbb{R} : x^2+x+1>0$''. Mệnh đề phủ định của mệnh đề đã cho là
	\choice{``$\forall x \in \mathbb{R} :x^2+x+1 < 0$''} 
		 {``$\forall x \in \mathbb{R} :x^2+x+1 = 0$''}
		{\True ``$\exists x \in \mathbb{R} :x^2+x+1  \le 0$''}
		{``$\exists x  \in \mathbb{R} :x^2+x+1 < 0$''}
	\loigiai{ Mệnh đề phủ định của mệnh đề ``$\forall x \in \mathbb{R} : x^2+x+1>0$'' là mệnh đề  ``$\exists x \in \mathbb{R} :x^2+x+1 \le 0$''.
	 }
\end{ex}

% nguồn: MaHoaNH_TracNghiem_Lop10
\begin{ex}%[0D1B1-5]%[ma:e96ca595]
	Cho  mệnh đề $P$: ``$\forall x \in \mathbb{R}, x^2+4 \ne 0$''. Mệnh đề $\overline{P}$ là
	\choice{$\overline{P}$: ``$\exists  x \in \mathbb{R}, x^2+4 \ne 0$''}
		{\True $\overline{P}$: ``$\exists x  \in \mathbb{R}, x^2+4 = 0$''}
		{$ \overline{P}$: ``$\forall x  \in \mathbb{R}, x^2+4 = 0$''}
		{$  \overline{P}$: ``$\forall x \in \mathbb{R}, x^2+4 > 0$''}
	\loigiai{Mệnh  đề phủ định của $P$ là $\overline{P}$: ``$\exists x \in \mathbb{R}, x^2+4 = 0$''.
	} 
\end{ex}

% nguồn: MaHoaNH_TracNghiem_Lop10
\begin{ex}%[0D1K1-5]%[ma:229507ee]
	Trong các mệnh đề sau, mệnh đề nào  \textbf{sai}?
	\choicew{0.4\textwidth} 
	\choice{$\forall n\in\mathbb{N}, n^2 \vdots 2\Rightarrow n \vdots 2$} 
		{$\forall n\in\mathbb{N}, n^2  \vdots 3\Rightarrow n \vdots 3$}
		{$\forall n\in\mathbb{N}, n^2 \vdots 6\Rightarrow  n \vdots 6$}
		 {\True $\forall n\in\mathbb{N}, n^2 \vdots 9\Rightarrow n \vdots 9$}
	\loigiai{Ta  thấy $3^2 \vdots 9$, nhưng $3\not\vdots 9$}
\end{ex}

% nguồn: MaHoaNH_TracNghiem_Lop10
\begin{ex}%[0D1B1-5]%[ma:6ed3495f]
	Phát biểu thành lời mệnh đề $"\exists x \in \mathbb{N} : x^2  +1 >10^{1000}"$.
	 \choice
		 {Tồn tại số nguyên $x$ sao cho $x^2 +1 >10^{1000}$}
		{Tồn tại số nguyên $x$ sao  cho $x^2+1 < 10^{1000}$}
		{Tồn tại số  nguyên dương $x$ sao cho $x^2 +1 >10^{1000}$}
		{\True Tồn tại  số tự nhiên $x$ sao cho $ x^2 +1 >10^{1000}$}
	 \loigiai{
	Cho mệnh đề chứa biến $P(x)$ với  $x\in X$. Khi đó, khẳng định: "Tồn tại $x$ thuộc $X$ để $P(x)$ đúng" là một mệnh đề. Kí hiệu là: "$"\exists x \in X: P(x)$".}
\end{ex}

% nguồn: MaHoaNH_TracNghiem_Lop10
\begin{ex}%[0D1K1-5]%[ma:577cd96b]
	Tìm mệnh  đề đúng.
	\choice 
		 {“$\exists x: \ x^2+3=0$”}
		 {“$\exists x: \ x^4+3x^2+2=0$”}
		 {“$\forall x\in \mathbb{Z}: \ x^5>x^2$”}
		{\True “$\forall n\in  \mathbb{N} \colon \left((2n+1)^2-1\right)\vdots 4$”}
	 \loigiai{
	$(2n+1)^2-1=4n^2+4n=4(n^2+n)\vdots 4, \forall n\in \mathbb{N}$. Vậy mệnh đề “$\forall n\in  \mathbb{N} \colon \left((2n+1)^2-1\right)\vdots 4$” đúng.}
\end{ex}

% nguồn: MaHoaNH_TracNghiem_Lop10
\begin{ex}%[0D1B1-5]%[ma:3a3a48ad]
	Tìm  mệnh đề sai.
	 \choice
		{“$\forall x:  \ x^2+2x+3>0$”}
		{\True “$\forall x:  \ x^2\geqslant x$”}
		{“$\exists x: \  x^2+5x+6=0$”}
		{“$\exists x: \ x<\dfrac{1}x$”} 
	 \loigiai{
	Chọn $x=\dfrac{1}{2}  \Rightarrow x^2<x$. Vậy mệnh đề “$\forall x: \ x^2\geqslant x$” sai.}
\end{ex}

% nguồn: MaHoaNH_TracNghiem_Lop10
\begin{ex}%[0D1B1-5]%[ma:91f38c8c]
	 Phủ định của mệnh đề $\forall n \in \mathbb{N},n^2-n$ là số chẵn?
	 \choice
		{$\forall n \in \mathbb{N},n^2-n$ là số  lẻ}
		 {$\forall n \in \mathbb{N},n^2-n$ là số chẵn}
		{$\exists n  \in \mathbb{N},n^2-n$ là số chẵn}
		{\True $\exists n  \in \mathbb{N},n^2-n$ là số lẻ}
	\loigiai{ 
	Mệnh đề phủ định của mệnh đề $"\forall x \in X:  P(x)"$ là $"\exists x \in X:\overline {P(x)}$".}
\end{ex}

% nguồn: MaHoaNH_TracNghiem_Lop10
\begin{ex}%[0D1B1-5]%[ma:00ca4101]
	Tìm mệnh đề phủ định của mệnh đề “$ \forall x\in \mathbb{R}\colon x^2>x$”. 
	\choice 
		 {$\forall x\in \mathbb{R}\colon x^2\leqslant x$}
		{$\exists  x\in \mathbb{R}\colon x^2>x$}
		{\True $\exists x\in  \mathbb{R}\colon x^2\leqslant x$}
		{$\exists x\in \mathbb{R}\colon  x^2<x$}
	\loigiai{ 
	Mệnh đề “$\forall x\in \mathbb{R}\colon x^2>x$” là mệnh đề “$\exists x\in \mathbb{R}\colon  x^2\leqslant x$”.}
\end{ex}

% nguồn: MaHoaNH_TracNghiem_Lop10
\begin{ex}%[0D1B1-5]%[ma:7dc496a8]
	Cho mệnh đề “Có một học sinh trong lớp C4  không chấp hành luật giao thông”. Mệnh đề phủ định của mệnh đề này là
	 \choice
		{Không có học sinh nào trong lớp C4  chấp hành luật giao thông}
		{\True Mọi học sinh trong lớp C4  đều chấp hành luật giao thông}
		{Có một học sinh trong lớp C4 chấp hành luật  giao thông}
		{Mọi học sinh trong  lớp C4 không chấp hành luật giao thông}
	 \loigiai{
	Mệnh đề phủ định là “ Mọi học sinh trong lớp C4 đều chấp  hành luật giao thông”.}
\end{ex}

% nguồn: MaHoaNH_TracNghiem_Lop10
\begin{ex}%[0D1B1-5]%[ma:b346572a]
	 Tìm mệnh đề phủ định của mệnh đề: $\forall x\in \mathbb{R}, \ x^2+x+5>0$.
	 \choice
		{$\exists  x\in \mathbb{R}, \ x^2+x+5<0$}
		{$\forall x\in \mathbb{R}, \ x^2+x+5<0$} 
		 {$\forall x\in \mathbb{R}, \ x^2+x+5\leqslant 0$}
		{\True $\exists x\in \mathbb{R}, \ x^2+x+5\leqslant  0$}
	 \loigiai{
	$\forall x\in \mathbb{R}, \ x^2+x+5>0$.  Suy ra mệnh đề phủ định là $\exists x\in \mathbb{R}, \ x^2+x+5\leqslant 0$.}
\end{ex}

% nguồn: MaHoaNH_TracNghiem_Lop10
\begin{ex}%[0D1K1-5]%[ma:e1e178ef]
	 Mệnh đề nào dưới đây \textbf{sai}?
	 \choice
		{$x(1-2x)\leqslant \dfrac{1}{8}, \ \forall x$} 
		{\True  $x^2+2+\dfrac{1}{x^2+2}>\dfrac{5}{2}, \ \forall x$}
		 {$\dfrac{x^2-x+1}{x^2+x+1}\geqslant \dfrac{1}{3}, \ \forall x$}
		{$\dfrac{x}{x^2+1}\leqslant \dfrac{1}{2}, \ \forall x$} 
	 \loigiai{
	Với  $x=0$ dễ thấy $x^2+2+\dfrac{1}{x^2+2}>\dfrac{5}{2}, \ \forall x$ sai.}
\end{ex}

% nguồn: MaHoaNH_TracNghiem_Lop10
\begin{ex}%[0D1B1-5]%[ma:89750980]
	Cho $x$ là số tự nhiên. Phủ định của mệnh đề  “$\forall x$ chẵn, $x^2+x$ là số chẵn” là mệnh đề:
	\choice 
		{$\exists x$ lẻ,  $x^2+x$ là số lẻ}
		{$\exists x$  lẻ, $x^2+x$ là số chẵn}
		{$\forall x$ lẻ, $x^2+x$  là số lẻ}
		{\True  $\exists x$ chẵn, $x^2+x$ là số lẻ}
	\loigiai{ 
	Mệnh đề phủ định là  “$\exists x$ chẵn, $x^2+x$ là số lẻ”.}
\end{ex}

% nguồn: MaHoaNH_TracNghiem_Lop10
\begin{ex}%[0D1B1-5]%[ma:a04503bd]
	Tìm mệnh đề phủ định của mệnh đề $"\forall z  \in \mathbb{Z} : z^2+z > z^4+10"$.
	\choice 
		{\True $"\exists z \in \mathbb{Z}  : z^2+z \leq z^4+10"$}
		{$"\exists z \in \mathbb{Z} : z^2+z  < z^4+10"$}
		{$"\exists z \in \mathbb{Z}  : z^2+z \geq z^4+10"$}
		{$"\exists z \in \mathbb{Z} : z^2+z >  z^4+10"$}
	\loigiai{ 
	 Mệnh đề phủ định của mệnh đề $"\forall x \in X: P(x)"$ là $"\exists x \in X:\overline {P(x)}$".}
\end{ex}

% nguồn: MaHoaNH_TracNghiem_Lop10
\begin{ex}%[0D1K1-5]%[ma:8d82f296]
	Tìm mệnh đề phủ định của  mệnh đề
	 $\forall m \in \mathbb{Z},\ \exists n \in \mathbb{Z} $: $m^2-n^2=1$.
	 \choice
		{\True $\exists m \in  \mathbb{Z},\ \forall n \in \mathbb{Z} $: $m^2-n^2 \neq 1$}
		{$\exists m \in \mathbb{Z},\ \forall n \in \mathbb{Z} $:  $m^2-n^2 = 1$}
		{$\exists m \in \mathbb{Z},\ \exists n \in \mathbb{Z} $: $m^2-n^2 \neq  1$}
		{$\forall m \in \mathbb{Z},\ \forall n \in \mathbb{Z}  $: $m^2-n^2 =1$}
	\loigiai{Mệnh đề phủ định của mệnh đề "$\forall m \in \mathbb{Z},\ \exists n \in \mathbb{Z} $: $m^2-n^2=1$" là mệnh đề "$\exists m  \in \mathbb{Z},\ \forall n \in \mathbb{Z} $: $m^2-n^2 \neq 1$".}
\end{ex}

% nguồn: MaHoaNH_TracNghiem_Lop10
\begin{ex}%[0D1K1-5]%[ma:e969adc9]
	 Tìm mệnh đề phủ định của mệnh đề
	 $\exists x \in \mathbb{R},\ \exists y \in \mathbb{R} $: $x^2-y^2 > 10^{1000}$.
	\choice 
		{$\exists x  \in \mathbb{R},\ \exists y \in \mathbb{R} $: $x^2-y^2 < 10^{1000}$}
		{$\forall x \in \mathbb{R},\ \forall y  \in \mathbb{R} $: $x^2-y^2 > 10^{1000}$}
		{$\forall x \in \mathbb{R},\ \forall y \in \mathbb{R} $: $x^2-y^2  < 10^{1000}$}
		{\True  $\forall x \in \mathbb{R},\ \forall y \in \mathbb{R} $: $x^2-y^2 \leq 10^{1000}$}
	\loigiai{Mệnh đề phủ  định của mệnh đề "$\exists x \in \mathbb{R},\ \exists y \in \mathbb{R} $: $x^2-y^2 > 10^{1000}$" là "$\forall x \in \mathbb{R},\ \forall y \in \mathbb{R} $: $x^2-y^2 \leq 10^{1000}$"}
\end{ex}

% nguồn: MaHoaNH_TracNghiem_Lop10
\begin{ex}%[0D1B1-5]%[ma:bc72cb64]
	Cho mệnh đề: “ Có một học sinh trong lớp 10A không thích học môn Toán”. Mệnh đề phủ định của mệnh đề này  là:
	\choice 
		{\True “ Mọi học sinh trong lớp 10A đều thích học môn  Toán”}
		{“ Mọi học sinh trong lớp 10A đều không thích học môn  Toán”}
		{“ Mọi học sinh  trong lớp 10A đều thích học môn Văn”}
		{“ Có  một học sinh trong lớp 10A thích học môn Toán”}
	 \loigiai{
	Phủ định của mệnh đề: “Có một học sinh trong lớp 10A không thích học môn Toán” là mệnh đề “Mọi học sinh  trong lớp 10A đều thích học môn Toán”.}
\end{ex}

% nguồn: MaHoaNH_TracNghiem_Lop10
\begin{ex}%[0D1B1-5]%[ma:a313fecc]
	Cho mệnh đề “$\forall x\in \mathbb{R}, \ x^2-x+7<0$”. Hỏi mệnh đề  nào là mệnh đề phủ định của mệnh đề trên?
	\choice 
		 {$\forall x\in \mathbb{R}, \ x^2-x+7>0$}
		{$ \nexists  x \in \mathbb{R}, \ x^2-x+7<0$}
		{\True $\exists x\in \mathbb{R}, \  x^2-x+7\geqslant 0$}
		{$\exists  x\in \mathbb{R}, \ x^2-x+7\leqslant 0$}
	 \loigiai{
	Phủ định của mệnh đề “$\forall x\in \mathbb{R}, \ x^2-x+7<0$” là mệnh đề  “$\exists x\in \mathbb{R}, \ x^2-x+7\geqslant 0$”.}
\end{ex}

% nguồn: MaHoaNH_TracNghiem_Lop10
\begin{ex}%[0D1B1-5]%[ma:24692917]
	Cho mệnh đề: “$\forall x\in \mathbb{R},\ x^2+3x+5>0$”. Mệnh đề phủ định của mệnh đề trên là 
	\choice 
		{$\forall  x\in \mathbb{R}, \ x^2+3x+5\leqslant 0$ }
		{\True $\exists x\in \mathbb{R},\ x^2+3x+5\leqslant  0$}
		{$\forall x\in \mathbb{R},  \ x^2+3x+5<0$ }
		{$\exists x\in \mathbb{R},  \ x^2+3x+5>0$ }
	 \loigiai{
	Chú ý: Phủ định của mệnh đề “$\forall x\in  \mathbb{R}\colon P(x)$ ” là “$\exists x\in \mathbb{R}\colon \overline{P(x)}$ ”.}
\end{ex}

% nguồn: MaHoaNH_TracNghiem_Lop10
\begin{ex}%[0D1K1-5]%[ma:c5c5ce8e]
	Mệnh đề \lq\lq Có ít nhất một số tự nhiên khác $0$\rq\rq\ mô tả  mệnh đề nào dưới đây?
	 \choice
		{\lq\lq$\forall n \in \mathbb{N}: n \neq  0$\rq\rq}
		{\lq\lq$\exists x \in \mathbb{N}:  x=0$\rq\rq}
		{\lq\lq$\exists x \in \mathbb{Z}:  x \neq 0$\rq\rq}
		{\True \lq\lq$\exists  x \in \mathbb{N}: x \neq 0$\rq\rq}
	\loigiai{ 
		Mô tả mệnh đề  \lq\lq$\exists x \in \mathbb{N}: x \neq 0$\rq\rq.
	 }
\end{ex}

% nguồn: MaHoaNH_TracNghiem_Lop10
\begin{ex}%[0D1K1-5]%[ma:f23d5e78]
	Mệnh  đề nào sau đây {\bf sai}?
	 \choice
		{$\forall x\in\mathbb{N}$, $x^2$ chia hết  cho $3$ $\Rightarrow$ $x$ chia hết cho $3$}
		{$\forall x\in\mathbb{N}$, $x$ chia hết cho $3$  $\Rightarrow$ $x^2$ chia hết cho $3$}
		{$\forall x\in\mathbb{N}$, $x^2$ chia hết cho $6$ $\Rightarrow$ $x$ chia hết cho $6$} 
		{\True $\forall x\in\mathbb{N}$, $x^2$  chia hết cho $9$ $\Rightarrow$ $x$ chia hết cho $9$}
	 \loigiai{
		\lq\lq $\forall x\in\mathbb{N}$, $x^2$ chia hết cho $9$ $\Rightarrow$ $x$ chia hết cho $9$\rq\rq\, sai vì chẳng hạn, chọn $x=3$,  khi đó $x^2=9$ chia hết cho $9$, còn $x=3$ không chia hết cho $9$.
	} 
\end{ex}

% nguồn: MaHoaNH_TracNghiem_Lop10
\begin{ex}%[0D1B1-5]%[ma:06f64290]
	Phủ định của mệnh đề $\exists x\in\mathbb{Z}\colon 1-x^2\geq 0$ là 
	 \choice
		{$\exists  x\in\mathbb{Z}\colon 1-x^2< 0$}
		 {$\forall x\in\mathbb{Z}\colon 1-x^2\geq 0$}
		{$\forall  x\in\mathbb{Z}\colon 1-x^2\neq 0$}
		 {\True $\forall x\in\mathbb{Z}\colon 1-x^2<0$}
	 \loigiai{
		Xét mệnh  đề $P\colon \exists x\in\mathbb{Z}\colon 1-x^2\geq 0 $ suy ra $ \overline{P}\colon \forall x\in\mathbb{Z}\colon 1-x^2<0$.
	 }
\end{ex}

% nguồn: MaHoaNH_TracNghiem_Lop10
\begin{ex}%[0D1B1-5]%[ma:eb67b49f]
	Cho $k$ là một số  nguyên tố. Gọi $\mathbb{P}$ là tập hợp các số nguyên tố, $\mathbb{I}$ là tập hợp các số vô tỉ. Mệnh đề $P$: "Với mọi số thực $x$, nếu $x$ là số vô tỉ thì $kx$ cũng là số vô tỉ" viết bằng các ký hiệu logic và tập hợp như thế nào?
	 \choice
		 {$P\colon k\in\mathbb{P}, \exists x\in\mathbb{R}; x\in\mathbb{I}\Rightarrow kx\in\mathbb{I}$}
		{\True $P\colon k\in\mathbb{P}, \forall x\in\mathbb{R};  x\in\mathbb{I}\Rightarrow kx\in\mathbb{I}$}
		{$P\colon \forall k\in\mathbb{P}, x\in\mathbb{R}, k\in\mathbb{P}; x\in\mathbb{I}\Leftrightarrow kx\in\mathbb{I}$} 
		 {$P\colon \forall k\in\mathbb{P}, x\in\mathbb{R}; kx\in\mathbb{I}\Rightarrow x\in\mathbb{I}$}
	 \loigiai{
	Mệnh đề $P$ viết bằng  các ký hiệu logic và tập hợp là $P\colon k\in\mathbb{P}, \forall x\in\mathbb{R}; x\in\mathbb{I}\Rightarrow kx\in\mathbb{I}$.
	 }
\end{ex}

% nguồn: MaHoaNH_TracNghiem_Lop10
\begin{ex}%[0D1B1-5]%[ma:136e451a]
	Xét mệnh đề $P: \lq\lq\forall x\in\mathbb{R}:  x^2\geq 0"$. Hãy xác định mệnh đề $\overline{P}.$
	\choice 
		 {\True $\overline{P}:\lq\lq\exists x\in\mathbb{R}: x^2<0"$}
		{$\overline{P}:\lq\lq\exists x\in\mathbb{R}:  x^2>0"$}
		 {$\overline{P}:\lq\lq\forall x\in\mathbb{R}: x^2<0"$}
		{$\overline{P}:\lq\lq\forall  x\in\mathbb{R}: x^2>0"$}
	 \loigiai{
		$P: \lq\lq\forall  x\in\mathbb{R}: x^2\geq 0"$ $\Rightarrow \overline{P}:\lq\lq\exists x\in\mathbb{R}: x^2<0".$
	 }
\end{ex}

% nguồn: MaHoaNH_TracNghiem_Lop10
\begin{ex}%[0D1B1-5]%[ma:ebfe5e66]
	Cho mệnh đề $P(x)\colon$ ``$\forall x\in\mathbb{R}\colon x^2+x+1>0$''. Phủ định của mệnh đề $P(x)$  là
	 \choice
		{ $\forall x\in\mathbb{R}\colon  x^2+x+1<0$}
		{$\forall x\in\mathbb{R}\colon  x^2+x+1\leq 0$}
		 {\True $\exists x\in\mathbb{R}\colon x^2+x+1\leq 0$}
		{$\exists x\in\mathbb{R}\colon x^2+x+1<  0$}
	\loigiai{ 
		Ta có $\overline{P(x)}\colon\exists x\in\mathbb{R}\colon x^2+x+1\leq  0$.
	 }
\end{ex}

% nguồn: MaHoaNH_TracNghiem_Lop10
\begin{ex}%[0D1B1-5]%[ma:a9e5d6d3]
	Mệnh đề phủ định  của mệnh đề “$\exists x\in \mathbb{R}, x^2+x+13=0$” là
	\choice 
		{\True “$\forall x\in \mathbb{R}, x^2+x+13\ne 0$”} 
		 {“$\exists x\in \mathbb{R}, x^2+x+13>0$”}
		{“$\forall x\in \mathbb{R}, x^2+x+13=0$”} 
		{“$\exists x\in \mathbb{R},  x^2+x+13\ne 0$”}
	\loigiai{ 
		Mệnh đề phủ định của mệnh đề “$\exists x\in \mathbb{R}, x^2+x+13=0$ ” là “$\forall  x\in \mathbb{R}, x^2+x+13\ne 0$”.
	} 
\end{ex}

% nguồn: MaHoaNH_TracNghiem_Lop10
\begin{ex}%[0D1B1-5]%[ma:57deedcf]
	Mệnh đề phủ định của mệnh  đề ``$\forall x\in \mathbb{R}, x^2-2x+1>0$'' là mệnh đề nào sau đây?
	 \choice
		{$\forall x\in  \mathbb{R}, x^2-2x+1<0$}
		 {$\exists x\in \mathbb{R}, x^2-2x+1>0$}
		{$\forall  x\in \mathbb{R}, x^2-2x+1\le 0$}
		 {\True $\exists x\in \mathbb{R}, x^2-2x+1\le 0$}
	\loigiai{ 
		Mệnh đề phủ định của mệnh đề ``$\forall x\in \mathbb{R}, x^2-2x+1>0$'' là mệnh đề $\exists  x\in \mathbb{R}, x^2-2x+1\le 0$.
	} 
\end{ex}

% nguồn: MaHoaNH_TracNghiem_Lop10
\begin{ex}%[0D1B1-5]%[ma:7e6cbfbb]
	Cho mệnh đề $A \colon$``$\forall x\in\mathbb{R}:x^2<x$''. Trong các mệnh đề sau, mệnh đề nào là phủ định của mệnh đề $A$? 
	\choice 
		{$\overline{A}  :$ ``$\exists x\in\mathbb{R}:x^2<x $''}
		{\True $\overline{A} :$ ``$\exists x\in\mathbb{R}:x^2\geq x$''} 
		{$\overline{A}  :$ ``$ \forall x\in\mathbb{R}:x^2>x$''}
		 {$\overline{A} :$ ``$\forall x\in\mathbb{R}:x^2\geq x$''}
	 \loigiai{
		Mệnh đề phủ định của mệnh đề $A \colon$``$\forall x\in\mathbb{R}:x^2<x$''  là $\overline{A} :$ ``$\exists x\in\mathbb{R}:x^2\geq x$''.
	} 
\end{ex}

% nguồn: MaHoaNH_TracNghiem_Lop10
\begin{ex}%[0D1B1-5]%[ma:7037a0cd]
	Tìm mệnh đề phủ định của mện đề  $P \colon$ \lq\lq $\forall x \in \mathbb{N}; \, x^2+x-1>0$ \rq\rq.
	 \choice
		{$\overline{P} \colon$  \lq\lq $\exists x \in \mathbb{N}; \, x^2 + x - 1>0$ \rq\rq}
		{$\overline{P} \colon$ \lq\lq  $\forall x \in \mathbb{N}; \, x^2 + x - 1>0$\rq\rq}
		{\True  $\overline{P} \colon$ \lq\lq $\exists x \in \mathbb{N}; \, x^2 + x - 1 \le 0$ \rq\rq}
		{$\overline{P} \colon$ \lq\lq$ \forall x \in \mathbb{N}; \, x^2  + x - 1 \le 0$ \rq\rq}
	\loigiai{ 
		Phủ định của mệnh đề $P \colon$ \lq\lq $\forall x \in \mathbb{N}; \, x^2+x-1>0$ \rq\rq \, là mệnh đề $\overline{P} \colon$ \lq\lq $\exists x \in  \mathbb{N}; \, x^2 + x - 1 \le 0$\rq\rq \, nên đáp án C đúng.
	 }
\end{ex}

% nguồn: MaHoaNH_TracNghiem_Lop10
\begin{ex}%[0D1B1-5]%[ma:c9f1dd40]
	Cho mệnh đề $P$: “$\left(2n+5\right)^2<  81$”. Mệnh đề phủ định của mệnh đề $P$ là
	 \choice
		{\True $\exists n \in \mathbb{N}$,  $\left(2n+5\right)^2\ge 81$}
		{$\forall n \in \mathbb{N}$, $\left(2n+5\right)^2\ge  81$}
		{$\exists n \in \mathbb{N}$, $\left(2n+5\right)^2\le 81$} 
		{$\exists  n \in \mathbb{N}$, $\left(2n+5\right)^2> 81$}
	\loigiai{ 
		Mệnh đề phủ định của mệnh đề $P$ là  "$\exists n \in \mathbb{N}$, $\left(2n+5\right)^2\ge 81$"
	} 
\end{ex}

% nguồn: MaHoaNH_TracNghiem_Lop10
\begin{ex}%[0D1B1-5]%[ma:e33f1294]
	Kí hiệu $X$ là tập  hợp các cầu thủ $x$ trong đội bóng đá. $P(x)$ là mệnh đề chứa biến ``$x$ cao trên $175$ (cm)''. Phát biểu thành lời mệnh đề ``$
	 \forall x\in X, P(x)$''?
	 \choice
		{\True Mọi cầu thủ trong  đội tuyển bóng đá đều cao trên $175$ (cm)}
		{Trong số các cầu thủ trong đội tuyển bóng đá có một số  cầu thủ cao trên $175$ (cm)}
		{Bất cứ ai cao trên $175$  (cm) đều là cầu thủ bóng đá}
		{Có một số người trong cao trên $175$ (cm) là  cầu thủ bóng đá}
	\loigiai{ 
		$\forall x\in X, P(x)\colon$ ``Mọi cầu thủ trong đội tuyển bóng đá đều cao  trên $175$ (cm)''.
	} 
\end{ex}

% nguồn: MaHoaNH_TracNghiem_Lop10
\begin{ex}%[0D1K1-5]%[ma:3887d523]
	Mệnh đề nào sau  đây là đúng?
	\choice 
		{$\exists  n \in \mathbb{Z},9n^2=1$}
		{$\forall n  \in \mathbb{N},n^2>n$}
		{$\exists  x \in \mathbb{Q},x^2-2=0$}
		{\True  $\exists y \in \mathbb{Z},3y^2-10y+3=0$}
	\loigiai{ 
		Ta  có $ 3y^2-10y+3=0\Leftrightarrow\hoac{&y=3\in\mathbb{Z}\\&y=\dfrac{1}{3}\not\in\mathbb{Z}.} $\\
		Vậy mệnh đề  $\exists y \in \mathbb{Z},3y^2-10y+3=0$ đúng.
	} 
\end{ex}

% nguồn: MaHoaNH_TracNghiem_Lop10
\begin{ex}%[0D1K1-5]%[ma:d67ade52]
	Mệnh  đề nào sau đây có mệnh đề phủ định đúng?
	\choice 
		{\True  ``$\forall n\in\mathbb{R}\colon 2n\geq n$''}
		 {``$\forall x\in\mathbb{R}\colon x<x+1$''}
		{``$\exists x\in\mathbb{Q}\colon  x^2=2$''}
		{``$\exists x\in\mathbb{R}\colon x^2-3x+1=0$''} 
	 \loigiai{
		Ta có mệnh đề ``$\forall n\in\mathbb{R}\colon 2n\geq n$'' sai, vì  với $ n=-1 $, $ 2n=-2<-1=n $.\\
		Suy ra mệnh đề phủ định là mệnh đề đúng. 
	} 
\end{ex}

% nguồn: MaHoaNH_TracNghiem_Lop10
\begin{ex}%[0D1B1-5]%[ma:bcc15775]
	Mệnh đề  nào sau đây là đúng?
	\choice 
		{\True $\forall x \in \mathbb{R}, \exists y \in  \mathbb{R},x+2y>3$}
		{$\forall x \in \mathbb{R},  \forall y \in \mathbb{R},x+2y>3$}
		{$\forall  x,y \in \mathbb{R},x^2+y^2+xy+x+y>0$}
		{$\exists m \in \mathbb{Z},m^2+1$ chia hết cho $4$} 
	 \loigiai{
		Mệnh đề $\forall x \in \mathbb{R}, \forall y  \in \mathbb{R},x+2y>3$ là mệnh đề sai, ví dụ cho $x=0, y=0$.\\
		Mệnh đề $\forall x,y \in \mathbb{R},x^2+y^2+xy+x+y>0$ là mệnh đề  sai, ví dụ cho $x=0, y=0$.\\
	Mệnh đề  $\exists m \in \mathbb{Z},m^2+1$ chia hết cho $4$ là mệnh đề sai. Vì nếu $m$ chẵn nghĩa là $m=2k$ thì $m^2+1=4k^2+1$ không chia hết cho $4$, nếu $m$ chẵn nghĩa là $m=2k+1$ thì $m^2+1=4k^2+4k+1$ không chia hết cho $4$.}
\end{ex}

% nguồn: MaHoaNH_TracNghiem_Lop10
\begin{ex}%[0D1B1-5]%[ma:c381cf96]
	Mệnh đề nào  sau đây là đúng?
	 \choice
		{\True $\exists x  \in \mathbb{Z},\dfrac{x^2}{x+2}\in\mathbb{Z}$}
		 {$\forall a,b \in \mathbb{R},a^2+b^2>2ab$}
		{$\exists x \in  \mathbb{R},x^2+3x+5=0$}
		 {$\forall y \in \mathbb{Z},y^3>y$}
	 \loigiai{
		Mệnh  đề $\exists x \in \mathbb{Z},\dfrac{x^2}{x+2}\in\mathbb{Z}$ là mệnh đề đúng, ví dụ cho $x=0$.\\
		Mệnh đề $\forall a,b \in \mathbb{R},a^2+b^2>2ab$ là mệnh đề sai, ví dụ cho $a=0,  b=0$.\\
		Mệnh đề $\exists  x \in \mathbb{R},x^2+3x+5=0$ là mệnh đề sai. Vì trên tập số thực phương trình $x^2+3x+5=0$ vô nghiệm.\\
	Mệnh đề $\forall y \in  \mathbb{Z},y^3>y$ là mệnh đề sai, ví dụ cho $y=-2$.}
\end{ex}

% nguồn: MaHoaNH_TracNghiem_Lop10
\begin{ex}%[0D1K1-5]%[ma:3123a644]
	Chọn mệnh  đề đúng trong các mệnh đề sau.
	 \choice
		{\True $\forall x \in  \mathbb{R}, x+1>x$}
		 {$\forall x \in \mathbb{R},|x|=x $}
		{$\exists  x \in \mathbb{R}, x-3 =x^2$}
		{$\exists x \in \mathbb{R},  x^2 <0$}
	\loigiai{ 
		Vì $x+1 >x \Leftrightarrow 1>0$ (luôn đúng $\forall x$) nên mệnh đề ở đáp án A đúng.\\ 
		Vì $\left|x\right| = x \Leftrightarrow x \geq 0$ nên mệnh  đề ở đáp án B sai.\\
		Vì $x-3=x^{2} \Leftrightarrow x^2 -x+3 =0 \Leftrightarrow \left(x-\dfrac{1}{2}\right)^2 +\dfrac{11}{4}=0$ (vô lý) nên mệnh đề ở đáp án C  sai.\\
		Vì $x^2 <0$ (vô lý)nên  mệnh đề ở đáp án D sai.\\
	} 
\end{ex}

% nguồn: MaHoaNH_TracNghiem_Lop10
\begin{ex}%[0D1B1-5]%[ma:75a3a7c5]
	Cho mệnh đề chứa biến $P(x)$: ``$x+15  \le x^2,\;x \in \mathbb{R}$''. Trong các mệnh đề sau, mệnh đề nào đúng?
	 \choice{$P(4)$}
		{\True  $P(5)$}
		 {$P(3)$}
		 {$P(0)$}
	\loigiai{Từ giả thiết ta có các kết quả  sau
		\begin{itemize}
			\item $P(4)$:  ``$4+15 \le 4^2$'' là mệnh đề sai.
			\item $P(5)$: ``$5+15 \le 5^2$'' là mệnh đề  đúng.
			\item $P(3)$:  ``$3+15 \le 3^2$'' là mệnh đề sai.
			\item $P(0)$: ``$0+15 \le 0^2$''  là mệnh đề sai.
		\end{itemize}
	} 
\end{ex}

% nguồn: MaHoaNH_TracNghiem_Lop10
\begin{ex}%[0D1K1-5]%[ma:2f7a9dc5]
	Mệnh đề nào sau đây là mệnh đề  đúng?
	\choice 
		{$\forall  x \in \mathbb{R},x>-2 \Rightarrow x^2>4$}
		{$\forall x \in  \mathbb{R},x^2>4 \Rightarrow x>2$}
		{\True $\forall x \in  \mathbb{R},x> 2 \Rightarrow x^2>4$}
		{$\forall x \in \mathbb{R},x^2>4 \Rightarrow x<-2$} 
	\loigiai{ 
		\begin{enumerate}[$\bullet$]
			\item Ta có $x=-1>-2, x^2<4$ nên  mệnh đề ``$\forall x \in \mathbb{R},x>-2 \Rightarrow x^2>4$'' sai.
			\item Ta có $(-5)^2>4$ nhưng  $-5<2$ nên mệnh đề ``$\forall x \in \mathbb{R},x^2>4 \Rightarrow x>2$'' sai.
			\item Ta có $\forall x \in  \mathbb{R}, x>2>0\Rightarrow x^2>4$.
			\item Ta có $5^2>4$  nhưng $5>2$ nên mệnh đề ``$\forall x \in \mathbb{R},x^2>4 \Rightarrow x<-2$'' sai.
		\end{enumerate}
	} 
\end{ex}

% nguồn: MaHoaNH_TracNghiem_Lop10
\begin{ex}%[0D1B1-5]%[ma:5d49f23d]
	 Mệnh đề nào sau đây đúng?
	 \choice
		{$\exists x\in\mathbb{R}: x^2<0$} 
		{$\forall  x\in\mathbb{R}: x>\dfrac{1}{x}$}
		{\True  $\forall x\in\left(-\infty;0\right]:|x|=-x$}
		{$\exists x\in\left(1;+\infty \right):  \sqrt{x-1}\leq 0$}
	 \loigiai{
		\begin{itemize}
			\item $\exists x\in\mathbb{R}: x^2<0$  sai vì $x^2\geq 0, \forall x\in \mathbb{R}$.
			\item $\forall x\in\mathbb{R}: x>\dfrac{1}{x}$ sai vì chẳng hạn khi chọn $x=\dfrac{1}{2}$ ta được $\dfrac{1}{2}>2$, vô lí. 
			\item $\forall x\in\left(-\infty;0\right]:|x|=-x$ đúng vì  $|x|=\heva{&x&\text{nếu } &x\geq 0\\&-x&\text{nếu } &x<0.}$
			\item $\exists x\in\left(1;+\infty \right): \sqrt{x-1}\leq  0$ sai vì $\sqrt{x-1}> 0,\forall x>1$.
		\end{itemize}
	 }
\end{ex}

% nguồn: MaHoaNH_TracNghiem_Lop10
\begin{ex}%[0D1B1-5]%[ma:45bded90]
	Trong các câu sau,  câu nào đúng?
	\choice 
		{Phủ định của mệnh đề \lq\lq$\forall  x\in \mathbb{Q}, 4x^2-1=0 $\rq\rq\ là mệnh đề \lq\lq$\forall x \in \mathbb{Q}, 4x^2-1>0$\rq\rq}
		{\True Phủ định của mệnh đề \lq\lq$\exists n \in \mathbb{N}, n^2+1 $  chia hết cho $4$\rq\rq\ là mệnh đề \lq\lq $\forall n \in \mathbb{N}, n^2+1$ không chia hết cho $4$\rq\rq}
		{Phủ định của mệnh đề \lq\lq$\forall x \in \mathbb{R}, (x-1)^2\ne x-1 $\rq\rq\ là mệnh đề \lq\lq$\forall x \in  \mathbb{R}, (x-1)^2=x-1$\rq\rq}
		{Phủ  định của mệnh đề \lq\lq$\forall n \in \mathbb{N}, n^2>n$\rq\rq\ là mệnh đề \lq\lq$\exists n\in \mathbb{N}, n^2<n$\rq\rq}
	\loigiai{ 
		\begin{itemize}
			\item Phủ định của mệnh đề \lq\lq$\forall x\in \mathbb{Q}, 4x^2-1=0 $\rq\rq\ là mệnh đề \lq\lq$\exists x \in  \mathbb{Q}, 4x^2-1\ne 0$\rq\rq.
			\item Phủ  định của mệnh đề \lq\lq$\forall x \in \mathbb{R}, (x-1)^2\ne x-1 $\rq\rq\ là mệnh đề \lq\lq$\exists x \in \mathbb{R}, (x-1)^2=x-1$\rq\rq.
			\item Phủ định của mệnh đề \lq\lq$\forall n \in  \mathbb{N}, n^2>n$\rq\rq\ là mệnh đề \lq\lq$\exists n\in \mathbb{N}, n^2\le n$\rq\rq.
		\end{itemize}
		Vậy phủ định của mệnh  đề \lq\lq$\exists n \in \mathbb{N}, n^2+1 $ chia hết cho $4$\rq\rq\ là mệnh đề \lq\lq $\forall n \in \mathbb{N}, n^2+1$ không chia hết cho $4$ \rq\rq\ là khẳng định đúng.
	} 
\end{ex}

% nguồn: MaHoaNH_TracNghiem_Lop10
\begin{ex}%[0D1K1-5]%[ma:2ea0d4db]
	Trong các mệnh đề sau, mệnh đề nào  đúng?
	\choice 
		{$\forall x\in  \mathbb{R}$, $(x-1)^2\ne x-1$}
		{$\forall x\in \mathbb{R}, |x|<3$ $  \Leftrightarrow x<3$}
		{$\exists n\in  \mathbb{N}, n^2+1$ chia hết cho $4$}
		{\True $\forall n\in \mathbb{N}, n^2+1$ không  chia hết cho $3$}
	\loigiai{ 
		\begin{itemize}
			\item Mệnh  đề ``$\forall x\in \mathbb{R}$, $(x-1)^2\ne x-1$'' sai vì với $x=1$ thì $(x-1)^2=x-1$.
			\item Mệnh đề ``$\forall x\in \mathbb{R},|x|<3$ $ \Leftrightarrow x<3$'' sai vì  khi $x=-4<3$ nhưng $|x|=4>3$.
			\item Mệnh đề ``$\exists n\in \mathbb{N},n^2+1$ chia hết cho  $4$'' sai vì\\
			+  Nếu $n=2k\left(k\in \mathbb{N}\right)$ thì $n^2+1=4k^2+1$ số này không chia hết cho $4$.\\
			+ Nếu $n=2k+1\left(k\in \mathbb{N}\right)$ thì $n^2+1=4k^2+4k+2$ số này cũng không chia hết cho  $4$.
			\item Mệnh đề ``$\forall n\in \mathbb{N},n^2+1$ không  chia hết cho $3$'' đúng vì\\
			+  Nếu $n=3k\left(k\in \mathbb{N}\right)$ thì $n^2+1=9k^2+1$ số này không chia hết cho $3$.\\
			+ Nếu $n=3k\pm  1\left(k\in \mathbb{N}^*\right)$ thì $n^2+1=9k^2\pm 6k+2$ số này không chia hết cho $3$.
		\end{itemize}
	 }
\end{ex}

% nguồn: MaHoaNH_TracNghiem_Lop10
\begin{ex}%[0D1K1-5]%[ma:2b89aaa2]
	Trong các mệnh đề sau,  mệnh đề nào \textbf{sai}?
	\choice 
		{$\exists n\in \mathbb{N}$, $n^2+11n+2$ chia  hết cho $11$}
		{\True $\exists n\in \mathbb{N}$, $n^2+1$ chia hết cho $4$} 
		{Tồn tại số nguyên  tố chia hết cho $5$}
		 {$\exists n\in \mathbb{Z}$, $2x^2-8=0$}
	 \loigiai{
		\begin{itemize}
			\item Xét  mệnh đề ``$\exists n\in \mathbb{N}$, $n^2+11n+2$ chia hết cho $11$''.\\
			Khi $n=3$ thì  giá trị của $(n^2+11n+2)$ bằng $44\vdots 11$ nên mệnh đề đúng.
			\item Xét mệnh đề ``$\exists n\in \mathbb{N}$, $n^2+1$  chia hết cho $4$''.\\
			Khi  $n=2k,k\in \mathbb{N} \Rightarrow n^2+1=4k^2+1$ không chia hết cho $4$, $k\in \mathbb{N}$.\\
			Khi $n=2k+1,k\in \mathbb{N} \Rightarrow n^2+1=(2k+1)^2+1=4k^2+4k+2$ không  chia hết cho $4$, $k\in \mathbb{N}$.
			\item Xét mệnh đề ``Tồn tại  số nguyên tố chia hết cho $5$''.\\
			Tồn tại số nguyên tố $5$ chia hết cho $5$  nên mệnh đề đúng.
			\item  Xét mệnh đề ``$\exists n\in \mathbb{Z}$, $2x^2-8=0$ ''.\\
			Phương trình $2x^2-8=0 \Leftrightarrow x^2=4 \Leftrightarrow x=-2;x=2\in  \Bbb{Z}$ nên mệnh đề đúng.
		\end{itemize}
	 }
\end{ex}


%%%%%%%%%%%%%%%% ĐÚNG/SAI %%%%%%%%%%%%%%%%

% nguồn: MaHoaNH_DungSai_Lop10
\begin{ex}%[0D1N1-5]%[ma:12445db2]
	Xét mệnh đề $P$: $\forall x\in\mathbb{R}$, \lq\lq$x^2+x+1>0$\rq\rq. Các câu sau là  đúng hay sai?
	 \choiceTF
		{Mệnh đề phủ định của mệnh đề $P$  là $P$: \lq\lq$\forall x\in\mathbb{R}$, $x^2+x+1>0$\rq\rq}
		{Mệnh đề phủ định của mệnh đề $P$ là $P$: \lq\lq$\exists x\in\mathbb{R}$,  $x^2+x+1>0$\rq\rq}
		{Mệnh đề phủ định của mệnh đề $P$  là $P$: \lq\lq$\exists x\in\mathbb{R}$, $x^2+x+1<0$\rq\rq}
		{\True Mệnh đề phủ định của mệnh đề  $P$ là mệnh đề sai}
	 \loigiai{
	\begin{itemchoice}
		\itemch Mệnh đề phủ định của mệnh đề $P$ là $\overline{P}$: \lq\lq$\exists x\in\mathbb{R}$,  $x^2+x+1\leq slant0$\rq\rq.
		\itemch Mệnh đề phủ định của mệnh đề $P$  là $\overline{P}$: \lq\lq$\exists x\in\mathbb{R}$, $x^2+x+1\leq slant0$\rq\rq.
		\itemch Mệnh đề  phủ định của mệnh đề $P$ là $\overline{P}$: \lq\lq$\exists x\in\mathbb{R}$, $x^2+x+1\leq slant0$\rq\rq.
		\itemch Với mọi $x\in\mathbb{R}$  ta có $x^2+x+1=\left(x+\dfrac{1}{2}\right)^2+\dfrac{3}{4}>0$. \\
		Do đó mệnh đề $P$  là mệnh đề đúng. \\
		Vậy mệnh đề phủ định của mệnh đề $P$ là mệnh đề  sai.
	\end{itemchoice}
	 }
\end{ex}

% nguồn: MaHoaNH_DungSai_Lop10
\begin{ex}%[0D1H1-5]%[ma:0161df68]
	 Các mệnh đề sau là đúng hay sai?
	\choiceTF 
		{$\forall x\in\mathbb{N}$, $\sqrt{x}$  là số vô tỉ}
		 {\True $\forall y\in\mathbb{R}$, $y^2\geq slant0$}
		 {\True $\exists n\in\mathbb{Z}$, $n+n=0$}
		{$\exists m\in\mathbb{Z}$,  $2m-1=0$}
	\loigiai{ 
	\begin{itemchoice}
		\itemch Với $x=4\in\mathbb{N}$, ta  có $\sqrt{x}=\sqrt{4}=2$ không là số vô tỉ.
		\itemch $\forall y\in\mathbb{R}$,  $y^2\geq slant0$.
		\itemch  Với $n=0\in\mathbb{Z}$, ta có $0+0=0$.
		\itemch Ta có $2m-1=0 \Leftrightarrow m=\dfrac{1}{2}\notin\mathbb{Z}$. 
	\end{itemchoice}
	 }
\end{ex}

% nguồn: MaHoaNH_DungSai_Lop10
\begin{ex}%[0D1H1-5]%[ma:e8579f0c]
	Xét tính đúng, sai của  mỗi mệnh đề sau.
	\choiceTF 
		{\True $\exists x \in \mathbb{Q},  4x^2-1=0$}
		{$\forall n  \in \mathbb{N}, n$ và $n+2$ là các số nguyên tố}
		{ $\forall x \in \mathbb{R},(x-1)^2  \neq x-1$}
		{  $\forall n \in \mathbb{N}, n^2>n$}
	\loigiai{ 
		\begin{itemchoice}
			\itemch Ta có $4x^2-1=0 \Leftrightarrow x=  \pm \dfrac{1}{2} \in \mathbb{Q}$.
			\itemch Ta cho $n=2 \in \mathbb{N}$ thì $n+2=4$  không là số nguyên tố.
			\itemch Ta cho  $x=1 \in \mathbb{R}$ thì $(x-1)^2=x-1=0$.
			\itemch Ta cho $n=0 \in \mathbb{N}$ thì $n^2=0$ nên  $n^2>n$ là sai.
		\end{itemchoice}
	 }
\end{ex}

% nguồn: MaHoaNH_DungSai_Lop10
\begin{ex}%[0D1H1-5]%[ma:e01e47de]
	Cho mệnh đề chứa biến $P(x):$ \lq\lq $x>\dfrac{1}{x}$ \rq\rq, xét tính  đúng sai của các mệnh đề sau:
	\choiceTF 
		 {$P(1)$}
		{\True  $P\left(-\dfrac{1}{3}\right)$}
		{ $\forall x \in \mathbb{N}, P(x)$} 
		{\True $\exists x \in \mathbb{N},  P(x)$}
	 \loigiai{
		\begin{itemchoice}
			\itemch Ta có $P(1)$ : \lq\lq $1>1$ \rq\rq\, đây là mệnh đề sai. 
			\itemch Ta có  $P\left(-\dfrac{1}{3}\right):-\dfrac{1}{3}>-3$ đây là mệnh đề đúng.
			\itemch Ta  có $\forall x \in \mathbb{N}, x>\dfrac{1}{x}$ là mệnh đề sai vì $P(1)$ là mệnh đề sai.
			 \itemch Ta có $\exists x \in \mathbb{N}, x>\dfrac{1}{x}$ là mệnh đề đúng vì $P(2): 2>\dfrac{1}{2}$ là mệnh đề đúng.
		\end{itemchoice}
	} 
\end{ex}

% nguồn: MaHoaNH_DungSai_Lop10
\begin{ex}%[0D1H1-5]%[ma:a8da277c]
	Xét tính  đúng sai của các mệnh đề sau:
	\choiceTF 
		{$A:$ \lq\lq $\forall x  \in \mathbb{R}, \forall y \in \mathbb{R}\colon x+y=1 $\rq\rq}
		{\True $B:$ \lq\lq $\exists x  \in \mathbb{R}, \exists y \in \mathbb{R}\colon x+y=2 $\rq\rq}
		{\True $C$ : \lq\lq $\exists x \in \mathbb{R}, \forall y \in \mathbb{R}\colon  y=x y$ \rq\rq}
		{\True $D$ : \lq\lq $\forall a \in \mathbb{R}, \exists b \in  \mathbb{R}\colon a=3 b$ \rq\rq}
	\loigiai{ 
		\begin{itemchoice}
			\itemch vì với $x=2,5; y=1 \Rightarrow P(2,5; 1)$ : \lq\lq $2,5 \in \mathbb{R},  1 \in \mathbb{R}\colon 2,5+1=1$ \rq\rq\, là mệnh đề sai.
			\itemch vì \lq\lq $1,5 \in \mathbb{R}, 0,5 \in \mathbb{R}\colon  1,5+0,5=2$ \rq\rq\, là mệnh đề đúng.
			\itemch vì $\exists x=1: y=1 . y$ đúng  $\forall y \in \mathbb{R}$.
			\itemch 
		\end{itemchoice}
	} 
\end{ex}

% nguồn: MaHoaNH_DungSai_Lop10
\begin{ex}%[0D1H1-5]%[ma:e7a88459]
	Xét tính đúng, sai  của mỗi mệnh đề sau.
	\choiceTF 
		{$\forall x \in  \mathbb{R}, x^2>0$}
		 {\True $\exists a \in \mathbb{Q}, a>a^2$}
		{\True $\forall n \in \mathbb{Z}, n^2+n+2$ chia hết  cho $2$}
		 { $\forall n \in \mathbb{N}, n(n+1)(n+2)$ không chia hết cho $3$}
	 \loigiai{
		\begin{itemchoice}
			\itemch Ta chọn $x=0$ thì  $x^2=0>0$ là sai.
			\itemch Ta chọn $a=\dfrac{1}{2} \in  \mathbb{Q}$ thì $a^2=\dfrac{1}{4}$ nên $a>a^2$ (đúng).
			\itemch Thật vậy: $\forall n  \in \mathbb{Z}, n^2+n+2=n(n+1)+2$, trong đó $n(n+1)$ là tích của hai số nguyên liên tiếp nên chia hết cho $2$ , vì vậy $n(n+1)+2$ cũng chia hết cho $2$.
			\itemch Ta cho $n=1$ thì $n(n+1)(n+2)=1 \cdot 2 \cdot 3=6$ chia hết cho  $3$.
		\end{itemchoice}
	} 
\end{ex}

% nguồn: MaHoaNH_DungSai_Lop10
\begin{ex}%[0D1H1-5]%[ma:09ced307]
	Xét tính đúng (sai)  của các mệnh đề sau
	 \choiceTF
		{$\forall x \in \mathbb{R},  x^3-x^2+1>0$}
		{\True $\exists n \in \mathbb{N},  n^2+3$ chia hết cho $4$}
		{ $P$ :\lq\lq $\forall x \in \mathbb{R}, \forall y \in  \mathbb{R}\colon x+y=1$ \rq\rq}
		{\True $Q:$  \lq\lq $\exists x \in \mathbb{R}, \exists y \in \mathbb{R}\colon x+y=2 $\rq\rq}
	 \loigiai{
		\begin{itemchoice}
			\itemch \lq\lq $\forall x \in  \mathbb{R}, x^3-x^2+10$ \rq\rq\, là mệnh đề sai
			\itemch \lq\lq $\exists n \in \mathbb{N}, n^2+3$  chia hết cho $4$ \rq\rq\, là mệnh đề đúng
			\itemch $P$ : \lq\lq $\forall x \in  \mathbb{R}, \forall y \in \mathbb{R}\colon x+y=1$ \rq\rq\, là mệnh đề sai
			Vì với $x=2{,}5; y=1 \Rightarrow P(2{,}5; 1): ~ 2{,}5 \in \mathbb{R}, 1 \in \mathbb{R}\colon 2{,}5+1=1$ \rq\rq là  mệnh đề sai.
			\itemch $Q$ : \lq\lq $\exists x \in \mathbb{R}, \exists  y \in \mathbb{R}\colon x+y=2$ \rq\rq\, là mệnh đề đúng
			Vì $Q(1{,}5; 0{,}5):$ ~ \rq\rq $1,5 \in \mathbb{R}, 0{,}5 \in \mathbb{R}\colon 1{,}5+0{,}5=2$ \rq\rq\,  là mệnh đề đúng.
		\end{itemchoice}
	 }
\end{ex}

% nguồn: MaHoaNH_DungSai_Lop10
\begin{ex}%[0D1H1-5]%[ma:783d1735]
	Xét tính  đúng, sai của các mệnh đề sau
	\choiceTF 
		{$\forall x\in  \mathbb{R}$, $x^3-x^2+1> 0$}
		{\True$\exists  n\in \mathbb{N}$, $n^2+3$ chia hết cho $4$}
		{$P\colon$ \lq\lq $\forall  x\in \mathbb{R}$, $\forall y\in \mathbb{R}\colon x+y=1$\rq\rq}
		{\True$Q\colon$\lq\lq $\exists x\in \mathbb{R}$, $\exists  y\in \mathbb{R}\colon x+y=2$\rq\rq}
	 \loigiai{
		\begin{itemchoice}
			\itemch \lq\lq $\forall x\in \mathbb{R}$, $x^3-x^2+1>0$\rq\rq ~là mệnh  đề sai.\\
			Vì với $ x=-1  $, ta có $ (-1)^3-(-1)^2+1=-1<0 $.
			\itemch \lq\lq $\exists n\in \mathbb{N}$,  $n^2+3$ chia hết cho $4$\rq\rq ~ là mệnh đề đúng.\\
			Vì có $ n=1 $ thì $ 1^2+3=4 $ chia hết cho  $ 4 $.
			\itemch $P\colon$ \lq\lq $\forall x\in \mathbb{R}$, $\forall  y\in \mathbb{R}\colon x+y=1$ \rq\rq ~ là mệnh đề sai.
			Vì với $x=2{,}5$; $y=1\Rightarrow  P(2{,}5;1)$ là mệnh đề sai.
			\itemch $Q\colon$\lq\lq $\exists x\in \mathbb{R},\exists y\in \mathbb{R}\colon x+y=2$\rq\rq ~ là mệnh đề  đúng.
			Vì $Q(1{,}5;0{,}5)$  là mệnh đề đúng.
		\end{itemchoice}
	} 
\end{ex}


%%%%%%%%%%%%%%%% TRẢ LỜI NGẮN %%%%%%%%%%%%%%%%

% nguồn: MaHoaNH_TraLoiNgan_Lop10
\begin{bt}%[0D1H1-5]%[ma:243251c2]
	Mệnh đề $\forall x \in \mathbb{R}, x^2-2+a>0$ với $a$ là số thực cho trước. Tìm giá  trị lớn nhất của $a$ để mệnh đề đúng.
	 \shortans{2}
	\loigiai{ 
		Vì $\forall x \in \mathbb{R}, x^2-2+a>0 \Leftrightarrow x^2>2-a \Leftrightarrow 2-a \leq 0 \Leftrightarrow a \leq  2$. Vậy $a=2$.
	} 
\end{bt}

% nguồn: MaHoaNH_TraLoiNgan_Lop10
\begin{bt}%[0D1H1-5]%[ma:5b6273b9]
	Mệnh đề \lq\lq$\forall x\in\mathbb{R},\, x^2-4x+a-2018 > 0$\rq\rq\, với a là số thực cho trước. Tìm giá trị nguyên lớn nhất của $a$ để mệnh đề đúng. 
	\shortans{2022} 
	 \loigiai{
		Ta có $x^2-4x+a-2\,018>0,\, \forall x\in\mathbb{R} \Leftrightarrow \Delta'<0 \Leftrightarrow 4-(a-2\,018)<0  \Leftrightarrow a>2\,022$.\\
	Vậy số  nguyên $a$ lớn nhất thỏa mãn là $a=2\,023$.}
\end{bt}

% nguồn: MaHoaNH_TraLoiNgan_Lop10
\begin{bt}%[0D1H1-5]%[ma:6c198254]
	Cho các mệnh đề sau: 
	\begin{enumerate} [a)]
		\item  $\sqrt{(-5)^2}=-5$.
		\item  $5^2+12^2=13^2$.
		 \item $\forall x \in \mathbb{R}, x^2+1>0$.
		\item $\exists x \in  \mathbb{Z}, x^2+3=0$.
	\end{enumerate}
	Trong các mệnh đề trên, có bao nhiêu mệnh  đề đúng?
	 \shortans{2}
	 \loigiai{
		\begin{enumerate} [a)]
			\item Mệnh đề sai vì $\sqrt{(-5)^2}=|-5|=5$. 
			\item Mệnh đề  đúng.
			\item Ta có $x^2 \geq 0$ và $1>0$ nên $\forall x \in \mathbb{R}, x^2+1>0$. Do đó mệnh đề  đúng.
			\item Giải phương trình $x^2+3=0$, phương trình vô nghiệm nên không có giá trị $x$ thỏa $x^2+3=0$ vì vậy mệnh  đề sai.
		\end{enumerate}
		Vậy ta có $2$  mệnh đề đúng.
	} 
\end{bt}

% nguồn: MaHoaNH_TraLoiNgan_Lop10
\begin{bt}%[0D1H1-5]%[ma:2aff37d4]
	 Trong các mệnh đề sau, có bao nhiêu mệnh đề đúng?
	\begin{enumerate} [a)]
		\item $A:$ \lq\lq $\forall x \in  \mathbb{R}, x^2 \geq 0$ \rq\rq.
		\item $B:$ \lq\lq Tồn tại  số tự nhiên đều là số nguyên tố\rq\rq.
		\item $C:$ \lq\lq $\exists  x \in \mathbb{N}, x$ chia hết cho $x+1$ \rq\rq.
		\item $D:$ \lq\lq $\forall n \in \mathbb{N}, n^4-n^2+1$  là hợp số \rq\rq.
		\item $E:$  \lq\lq Tồn tại hình thang là hình vuông \rq\rq.
		\item $F:$ \lq\lq  Tồn tại số thực $a$ sao cho $a+1+\dfrac{1}{a+1} \leq 2$ \rq\rq.
	\end{enumerate}
	 \shortans{5}
	\loigiai{ 
		\begin{enumerate} [a)]
			\item Mệnh đề $A$ đúng và $\overline{A}: \exists x \in R, x^2<0$. 
			\item Mệnh đề $B$ đúng và $\overline{B}$ : \lq\lq Với mọi số tự  nhiêu đều không phải là số nguyên tố\rq\rq.
			\item Mệnh đề $C$ đúng vì cho $x=0$ và $\overline{C}:$  ~ \rq\rq $\forall x \in \mathbb{N}, x \mid(x+1) $\rq\rq
			\item Mệnh đề  $D$ sai vì với $n=2$ ta có $n^4-n^2+1=13$ không phải là hợp số.\\
			Mệnh đề phủ định là $\overline{D}$ : \lq\lq $\exists n \in N, n^4-n^2+1$ là số số  nguyên tố\rq\rq.
			\item Mệnh đề $E$ đúng và $\overline{E}$ : \lq\lq Với mọi hình thang đều không là hình vuông \rq\rq. 
			\item Mệnh đề $F$ đúng và  mệnh đề phủ định là $\overline{F}$ \lq\lq Với mọi số thực $a$ thì $a+1+\dfrac{1}{a+1}>2$ \rq\rq.
		\end{enumerate}
	} 
\end{bt}

% nguồn: MaHoaNH_TraLoiNgan_Lop10
\begin{bt}%[0D1V1-5]%[ma:ed38142b]
	Mệnh đề \lq\lq$ \forall x\in\mathbb{R},\, x^2+5-m\ge 0 $\rq\rq~với $m$ là số thực cho trước. Tìm giá trị lớn nhất của $m$ để  được mệnh đề đúng.
	\shortans{5} 
	 \loigiai{
		Ta có $\forall x\in\mathbb{R},\, x^2+5-m\ge 0\Leftrightarrow \forall x\in\mathbb{R},\, x^2\ge m-5  \Leftrightarrow m-5\le 0\Leftrightarrow m\le 5$.\\
	Vậy giá trị lớn nhất của $m$  thỏa đề là: $m=5$.}
\end{bt}

% nguồn: MaHoaNH_TraLoiNgan_Lop10
\begin{bt}%[0D1V1-5]%[ma:3bb6a4ce]
	Xét mệnh đề \lq\lq $\forall x \in \mathbb{R}$, $x^2+2-a>0$\rq\rq với $a$ là một số thực cho trước. Có tất  cả bao nhiêu giá trị nguyên dương của $a$ để mệnh đề đang xét là mệnh đề đúng>
	\shortans{1} 
	 \loigiai{
		Vì $x^2+2-a>0$, $\forall x \in  \mathbb{R} \Leftrightarrow x^2>a-2$, $\forall x \in \mathbb{R}$.\\
		Mà $x^2 \geq 0$, $x \in \mathbb{R}$, do đó $a-2<0 \Leftrightarrow a<2$. 
	Vì $a$ nguyên dương  nên có $ 1 $ giá trị thỏa mãn.}
\end{bt}


%%%%%%%%%%%%%%%% TỰ LUẬN %%%%%%%%%%%%%%%%

% nguồn: MaHoaNH_TuLuan_Lop10
\begin{bt}%[0D1H1-5]%[ma:1dd675fc]
	Phát biểu bằng lời mệnh đề \lq\lq $\forall x \in \mathbb{R}, x^2+1 \leq 0$\rq\rq \, và cho biết mệnh đề đó  đúng hay sai?
	 \loigiai{
		Mệnh đề phát biểu là \lq\lq Mọi số thực đều có bình phương cộng với $1$ nhỏ  hơn hoặc bằng $0$ \rq\rq.\\
		Mệnh  đề sai do $x^2 \geq 0$ nên $x^2+1>0$.
	 }
\end{bt}

% nguồn: MaHoaNH_TuLuan_Lop10
\begin{bt}%[0D1H1-5]%[ma:a3794ebc]
	 Viết mệnh đề phủ định của mệnh đề $A\colon$ \lq\lq $\forall x \in \mathbb{R},-4x^2+4x-1 \leq 0$ \rq\rq \, và xét tính đúng, sai của mệnh đề đó.
	\loigiai{ 
		Ta có $A:$ \lq\lq$\forall  x \in \mathbb{R},-4x^2+4x-1<0$\rq\rq\, là mệnh đề sai vì\\ $-4x^2+4x-1<0 \Leftrightarrow-(2x-1)^2<0 \Leftrightarrow x \neq \dfrac{1}{2}$.\\
		Khi đó mệnh đề phủ định \lq\lq$ \exists x  \in \mathbb{R},\,-4x^2+4x-1>0$\rq\rq \,\, là mệnh đề đúng.
	} 
\end{bt}

% nguồn: MaHoaNH_TuLuan_Lop10
\begin{bt}%[0D1H1-5]%[ma:56f5b7d3]
	Xét tính đúng sai và viết mệnh đề phủ  định của các mệnh đề sau:
	\begin{enumerate} [a)]
		\item $\forall  x \in \mathbb{R}, x^2+2x+2>0$.
		\item $\exists  x \in \mathbb{R}, x^2+3x+4=0$.
	\end{enumerate}
	\loigiai{ 
		\begin{enumerate} [a)]
			\item Mệnh đề đúng, vì $x^2+2x+2=x^2+2x+1+1=(x+1)^2+1>0$ với mọi  số thực $x$.\\
			Mệnh đề phủ  định của mệnh đề này là $\exists x \in \mathbb{R}, x^2+2x+2 \leq 0$.
			\item Mệnh đề sai, vì phương trình $x^2+3x+4=0$ vô nghiệm $(\Delta=-7<0)$.\\ 
			Mệnh đề phủ định của mệnh đề này là $\forall x \in  \mathbb{R}, x^2+3x+4 \neq 0$.
		\end{enumerate}
	} 
\end{bt}

% nguồn: MaHoaNH_TuLuan_Lop10
\begin{bt}%[0D1H1-5]%[ma:0099f11d]
	Lập mệnh đề phủ định của mỗi mệnh đề  sau và xét tính đúng sai của mỗi mệnh đề phủ định đó
	\begin{listEX}[2]
		\item $\forall  x \in \mathbb{R}$, $x^2 \neq 2x-2$.
		\item  $\forall x \in \mathbb{R}$, $x^2 \leq 2x-1$.
	\end{listEX}
	\loigiai{ 
		\begin{enumerate}
			 \item $\exists x \in \mathbb{R}$, $x^2=2x-2$.\\
			Mệnh  đề này sai vì phương trình $x^2-2x+2=0$ vô nghiệm trên tập số thực.
			\item $\exists x \in \mathbb{R}$, $x^2  > 2x-1$.\\
			Mệnh đề này đúng vì với $x=2$ thì $2^2>2\cdot 2-1$. 
		\end{enumerate}
	} 
\end{bt}

% nguồn: MaHoaNH_TuLuan_Lop10
\begin{bt}%[0D1H1-5]%[ma:6a1f5e36]
	Dùng kí hiệu $\forall$ hoặc $\exists$ để mô tả các mệnh đề  sau:
	\begin{enumerate} [a)]
		\item Với mọi số tự nhiên $x$,  $\sqrt{x}$ là số vô tỉ.
		\item Tồn tại số  nguyên cộng với chính nó bằng $0$.
		\item  Bình phương của mọi số thực đều không âm.
		\item Có số hữu tỉ $n$  sao cho $3 n-1=0$.
	\end{enumerate}
	 \loigiai{
		\begin{enumerate} [a)]
			\item $\forall  x \in \mathbb{N}$, $\sqrt{x}$ là số vô tỉ.
			\item  $\exists x \in \mathbb{Z}$, $x+x=0$.
			 \item $\forall x \in \mathbb{R}$, $x^2 \geq 0$.
			\item $\exists n \in  \mathbb{Q}$, $3n-1=0$.
		\end{enumerate}
	 }
\end{bt}

% nguồn: MaHoaNH_TuLuan_Lop10
\begin{bt}%[0D1H1-5]%[ma:2e10abda]
	Sử dụng kí hiệu \lq\lq$\forall$\rq\rq \,để viết mỗi mệnh đề sau và xét xem mệnh đề đó là  đúng hay sai, giải thích vì sao.
	\begin{enumerate}
		\item $P\colon$\lq\lq Với mọi số  thực $x$, $x^2+1>0$\rq\rq.
		\item $Q\colon$\lq\lq Với mọi số tự nhiên $n$, $n^2+n$ chia  hết cho $6$\rq\rq.
	\end{enumerate}
	 \loigiai{
		\begin{enumerate}
			\item $P\colon$\lq\lq Với mọi số  thực $x$, $x^2+1>0$\rq\rq.\\
			Mệnh đề  được viết là $P \colon \lq\lq\forall x \in \mathbb{R}$, $x^2+1>0$\rq\rq.\\
			Xét một số thực $x$  tùy ý, ta phải chứng tỏ rằng $x^2+1>0$.\\
			Thật vậy,  ta có $x^2+1 \geq 1>0$.\\
			Vậy mệnh đề $P$  là mệnh đề đúng.
			\item $Q\colon$\lq\lq Với mọi  số tự nhiên $n$, $n^2+n$ chia hết cho $6$\rq\rq.\\
			Mệnh đề được viết là $Q\colon\lq\lq \forall n \in  \mathbb{N}$, $\left(n^2+n\right) \,\vdots\, 6$\rq\rq.\\
			Để chứng minh mệnh đề $Q$ là sai,  ta cần chỉ ra một giá trị cụ thể của $n$ để nhận được mệnh đề sai.\\
			Thật vậy, chọn $n=1$, ta thấy $n^2+n=2$ không chia hết cho $6$.\\ 
			Vậy mệnh đề  $Q$ là mệnh đề sai.
		\end{enumerate}
	 }
\end{bt}

% nguồn: MaHoaNH_TuLuan_Lop10
\begin{bt}%[0D1B1-5]%[ma:aaf3c4ee]
	Viết lại mệnh đề  \textit{$P\colon$``Mọi số thực đều nhỏ hơn nghịch đảo của nó''} bằng cách sử dụng kí hiệu $\forall$ hoặc $\exists$, xét tính đúng sai và nêu mệnh đề phủ định của nó .
	\loigiai{ 
		$P\colon$ \lq\lq$\forall  x \in \mathbb{R}, x < \dfrac{1}{x}$\rq\rq. \quad (Mệnh đề sai)\\
		$\overline{P}\colon$ \lq\lq$\exists x \in \mathbb{R}, x  \ge \dfrac{1}{x}$\rq\rq.
	 }
\end{bt}

% nguồn: MaHoaNH_TuLuan_Lop10
\begin{bt}%[0D1B1-5]%[ma:ff4d96fc]
	Cho mệnh đề chứa biến $P(x):\text{\lq\lq}x^2=x\text{\rq\rq}$. Xác định tính đúng - sai của  các mệnh đề sau: $P(0)$; $P(-1)$; $P(1)$; \lq\lq$\exists x\in \mathbb{R}$, $P(x)$\rq\rq; \lq\lq$\forall x\in \mathbb{R}$, $\overline{P(x)}$\rq\rq.
	 \loigiai{
		\begin{itemize}
			\item  $P(0):\text{\lq\lq}0^2=0\text{\rq\rq}$ đúng.
			\item $P(-1):\text{\lq\lq}(-1)^2=-1\text{\rq\rq}$  sai.
			\item  $P(1):\text{\lq\lq}1^2=1\text{\rq\rq}$ đúng.
			 \item \lq\lq$\exists x\in \mathbb{R}$, $P(x)$\rq\rq\ hay \lq\lq$\exists x\in \mathbb{R}$, $x^2=x$\rq\rq\ đúng vì $x=0$ thỏa mãn.
			\item \lq\lq$\forall x\in \mathbb{R}$, $\overline{P(x)}$\rq\rq\ hay \lq\lq$\forall  x\in \mathbb{R}$, $x^2\ne x$\rq\rq\ sai vì $x=0$ không thỏa mãn.
		\end{itemize}
	} 
\end{bt}

% nguồn: MaHoaNH_TuLuan_Lop10
\begin{bt}%[0D1B1-5]%[ma:e1a04003]
	Cho mệnh đề chứa biến $P(x):\text{\lq\lq}x^3-2x=0\text{\rq\rq}$. Xác định tính đúng - sai của các mệnh đề sau: $P(0)$; $P(2)$; $P(\sqrt{2})$; \lq\lq$\exists x\in \mathbb{R}$, $P(x)$\rq\rq; \lq\lq$\forall x\in \mathbb{R}$, $\overline{P(x)}$\rq\rq. 
	\loigiai{ 
		\begin{itemize}
			 \item $P(0):\text{\lq\lq}(0)^3-2(0)=0\text{\rq\rq}$ đúng.
			 \item $P(2):\text{\lq\lq}(2)^3-2(2)=0\text{\rq\rq}$ sai.
			 \item $P(\sqrt{2}):\text{\lq\lq}(\sqrt{2})^3-2(\sqrt{2})=0\text{\rq\rq}$ đúng.
			\item \lq\lq$\exists x\in \mathbb{R}$,  $P(x)$\rq\rq\ hay \lq\lq$\exists x\in \mathbb{R}$, $x^3-2x=0$\rq\rq\ đúng.
			\item \lq\lq$\forall x\in \mathbb{R}$, $\overline{P(x)}$\rq\rq\ hay \lq\lq$\forall x\in  \mathbb{R}$, $x^3-2x\ne 0$\rq\rq\ sai.
		\end{itemize}
	} 
\end{bt}

% nguồn: MaHoaNH_TuLuan_Lop10
\begin{bt}%[0D1K1-5]%[ma:022b1fdf]
	Các mệnh đề sau đúng hay sai?  Vì sao? Viết mệnh đề phủ định của chúng.
	\begin{enumerate}
		\item  \lq\lq $\exists x \in \mathbb{N},~ x~\text{chia hết cho}~x+1$\rq\rq.
		\item \lq\lq $\forall x \in \mathbb{Z},~ x \geq -1 \Rightarrow x^2 \geq  1$\rq\rq,
	\end{enumerate}
	 \loigiai{
		\begin{enumerate}
			\item Mệnh đề đã cho đúng vì tồn tại $x=0 \in \mathbb{N}$ thì $0$ chia hết  cho $1$ là điều đúng. Nó có mệnh đề phủ định là \lq\lq $\forall x \in \mathbb{N},~ x~\text{không chia hết cho}~x+1$\rq\rq.
			\item Mệnh đề đã cho sai vì tồn tại $x=0 \in \mathbb{Z}$ thì \lq\lq $0 \geq -1 \Rightarrow 0^2 \geq 1$\rq\rq\ là điều sai.\\ 
			Áp dụng kiến thức $\overline{A \Rightarrow B} = A \wedge \overline{B}$ với $A \wedge \overline{B}$ là giao của hai mệnh đề $A$ và  $\overline{B}$ ta có mệnh đề phủ định của mệnh đề đã cho là \lq\lq $\exists x \in \mathbb{Z},~ x \geq -1 ~\text{và}~ x^2 \leq 1$\rq\rq,
		\end{enumerate}
	 }
\end{bt}

% nguồn: MaHoaNH_TuLuan_Lop10
\begin{bt}%[0D1B1-5]%[0D1B4-2]%[ma:7cc11452]
	\begin{enumerate}
		\item Cho mệnh đề $A:$ \lq \lq $\exists \in \mathbb{R} \colon x^2-2x+7=0$\rq  \rq. Tìm phủ định của $A$.
		\item  Cho tập hợp $A=\{ x\in \mathbb{R} | -1 < x \le 4 \}$ và $B=\{ x\in \mathbb{R} | x \le 1 \}.$
		Tìm  các tập hợp: $A \cap B$, $A \cup B$, $B \setminus A$, $C_{\mathbb{R}}B.$
	\end{enumerate}
	\loigiai{ 
		\begin{enumerate}
			\item Ta có  $\overline{A}:$ \lq \lq$\forall x \in \mathbb{R}:x^2-2x+7\ne 0$\rq \rq.
			\item Ta  có $A=\left(-1;4\right]$ và $B=\left(- \infty;1 \right]$ nên:\\
			$A \cap B = \left(-1;1  \right].$\\
			$A \cup B = \left(- \infty;  4 \right].$\\
			$B \setminus A = \left(- \infty; -1 \right].$\\ 
			 $C_{\mathbb{R}}B = \left(1; + \infty \right).$
		\end{enumerate}
	 }
\end{bt}

% nguồn: MaHoaNH_TuLuan_Lop10
\begin{bt}%[0D1K1-5]%[ma:6f357c44]
	Xác định tính đúng - sai  của các mệnh đề sau:
	\begin{multicols}{2}
		\begin{enumerate}[a.]
			\item $\forall x\in\mathbb{R}: x>-2\Rightarrow  x^2>4$.
			\item $\forall x\in\mathbb{R}: x>2\Rightarrow x^2>4$. 
			\item $\forall m,n\in\mathbb{N}: m$ và $n$ là  các số lẻ $\Leftrightarrow m^2 +n^2$ là số chẵn.
			\item  $\forall x\in\mathbb{R}: x^2>4\Rightarrow x>2$.
		\end{enumerate}
	\end{multicols}
	 \loigiai{
		\begin{enumerate}[a.]
			\item $\forall x\in\mathbb{R}: x>-2\Rightarrow x^2>4$.\\Đây là mệnh đề sai, vì $x=-1$ thì \lq\lq$-1>-2\Rightarrow (-1)^2>4$\rq\rq\,  là mệnh đề sai.
			\item $\forall x\in\mathbb{R}: x>2\Rightarrow x^2>4$.\\Đây là mệnh đề đúng, vì với mọi số thực lớn hơn $2$ thì bình phương của nó luôn lớn hơn  $4$.
			\item $\forall m,n\in\mathbb{N}: m$ và $n$ là các số lẻ $\Leftrightarrow m^2 +n^2$ là số chẵn.\\Đây là mệnh đề  sai, vì \lq\lq $m=2k+1$, $n=2l+1$ với $k,l\in\mathbb{N}$ thì $m^2+n^2$ là số chẵn\rq\rq\,là đúng. Tuy nhiên \lq\lq $m^2+n^2$ là số chẵn thì $m$, $n$ là số lẻ\rq\rq\, là sai. Do đó mệnh đề tương đương này sai.
			\item $\forall x\in\mathbb{R}: x^2>4\Rightarrow x>2$.\\Đây là mệnh đề sai,  vì $x=-3$ thì \lq\lq$(-3)^2>4\Rightarrow -3>2$\rq\rq\, là mệnh đề sai.
		\end{enumerate}
	} 
\end{bt}

% nguồn: MaHoaNH_TuLuan_Lop10
\begin{bt}%[0D1Y1-5]%[ma:39aa458d]
	Sử dụng kí hiệu \lq\lq$\exists$\rq\rq\, để viết mỗi mệnh đề sau và xét  xem mệnh đề đó là đúng hay sai? Giải thích vì sao?
	\begin{enumerate}
		\item  $M\colon$\lq\lq Có ít nhất một số thực $x$ sao cho $x^3=-8$\rq\rq.
		\item $N\colon$\lq\lq Tồn tại số nguyên $x$ sao cho  $2x+1=0$\rq\rq.
	\end{enumerate}
	 \loigiai{
		\begin{enumerate}
			\item $M\colon$\lq\lq Có ít nhất một  số thực $x$ sao cho $x^3=-8$\rq\rq.\\
			Mệnh đề được viết là  $M\colon\lq\lq\exists x \in \mathbb{R}, x^3=-8$\rq\rq.
			Để chứng tỏ mệnh đề $M$ là đúng, ta cần chỉ ra một giá  trị cụ thể của $x$ để nhận được mệnh đề đúng.\\
			 Thật vậy, chọn $x=-2$, ta thấy $(-2)^3=-8$.\\
			Vậy mệnh đề $M$ là mệnh đề đúng.\\ 
			 Mệnh đề $N\colon\lq\lq\exists x \in \mathbb{Z}, 2x+1=0$\rq\rq.
			\item $N\colon$\lq\lq Tồn  tại số nguyên $x$ sao cho $2x+1=0$\rq\rq.\\
			Để chứng minh mệnh đề $N$ là sai, ta phải chứng tỏ  rằng với số nguyên $x$ tùy ý thì $2x+1 \neq 0$.\\
			Thật vậy,  xét một số nguyên $x$ tùy ý, ta có $2x+1 \neq 0$.\\
			Vì thế mệnh  đề $N$ là mệnh đề sai.
		\end{enumerate}
	 }
\end{bt}

% nguồn: MaHoaNH_TuLuan_Lop10
\begin{bt}%[0D1B1-5]%[ma:2fc0a578]
	Các mệnh đề sau đúng hay  sai? Nếu sai hãy sửa lại để có một mệnh đề đúng.
	\begin{listEX}[2]
		\item $\forall  x\in \mathbb{R}$, $x^2>x$.
		 \item $\forall x\in \mathbb{R}$, $x^2+y^2\le 2xy$.
		\item $\exists x\in \mathbb{N}$,  $x^2\le x$.
		\item $\exists n\in \mathbb{N}$,  $n^2+n+1\vdots 7$.
	\end{listEX}
	\loigiai{ 
		\begin{listEX}
			 \item $\forall x\in \mathbb{R}$, $x^2>x$.\\
			Mệnh đề sai. Sửa lại thành \lq\lq $\exists  x\in \mathbb{R}$, $x^2>x$\rq\rq.
			\item $\forall x\in \mathbb{R}$, $x^2+y^2\le 2xy$.\\ 
			Mệnh  đề sai. Sửa lại thành \lq\lq $\forall x\in \mathbb{R}$, $x^2+y^2\ge 2xy$\rq\rq.
			\item $\exists x\in \mathbb{N}$, $x^2\le x$.\\ 
			Mệnh đề  đúng.
			 \item $\exists n\in \mathbb{N}$, $n^2+n+1\vdots 7$.\\
			 Mệnh đề đúng.
		\end{listEX}
	} 
\end{bt}

% nguồn: MaHoaNH_TuLuan_Lop10
\begin{bt}%[0D1B1-5]%[ma:4a8ba17d]
	Phát biểu thành lời các mệnh đề sau và xét tính đúng sai của các mệnh đề đó. 
	\begin{listEX}[2]
		 \item $\forall x\in \mathbb{R}$, $x^2>0$.
		 \item $\exists n\in \mathbb{N}$, $n^2=n$.
		\item $\exists n\in  \mathbb{N}$, $n\le 2n$.
		\item  $\exists x\in \mathbb{R}$, $x<0$.
		\item $\forall  x\in \mathbb{N}$, $1{,}2<|x|<2{,}1$.
		 \item $\forall n\in \mathbb{N}$, $n^2+1$ chia hết cho $3$.
	\end{listEX}
	\loigiai{ 
		\begin{listEX}
			\item $\forall x\in  \mathbb{R}$, $x^2>0$.\\
			Phát biểu \lq\lq Bình phương của một số thực đều lớn hơn $0$\rq\rq.  Mệnh đề sai vì $x=0$ không thỏa mãn.
			\item  $\exists n\in \mathbb{N}$, $n^2=n$.\\
			Phát biểu \lq\lq Tồn tại số tự nhiên bằng bình phương của chính nó\rq\rq.  Mệnh đề đúng vì $n=1$ thỏa mãn.
			\item $\exists n\in  \mathbb{N}$, $n\le 2n$.\\
			Phát biểu \lq\lq Tồn tại số tự nhiên bé hơn  hoặc bằng hai lần số đó\rq\rq. Mệnh đề đúng vì $n=1$ thỏa mãn.
			 \item $\exists x\in \mathbb{R}$, $x<0$.\\
			Phát biểu \lq\lq Tồn  tại số thực bé hơn $0$\rq\rq. Mệnh đề đúng vì $x=-1$ thỏa mãn.
			\item $\forall  x\in \mathbb{N}$, $1{,}2<|x|<2{,}1$.\\
			Phát biểu \lq\lq Giá trị  tuyệt đối của tất cả các số tự nhiên đều lớn hơn $1{,}2$ và bé hơn $2{,}1$\rq\rq. Mệnh đề sai vì $x=3$ không thỏa mãn.
			\item  $\forall n\in \mathbb{N}$, $n^2+1$ chia hết cho $3$.\\
			Phát  biểu \lq\lq Tất cả số tự nhiên $n$ đều thỏa mãn $n^2+1$ chia hết cho $3$\rq\rq. Mệnh đề sai vì $n=1$ không thỏa mãn.
		\end{listEX}
	} 
\end{bt}

% nguồn: MaHoaNH_TuLuan_Lop10
\begin{bt}%[0D1B1-5]%[ma:b5c59cdf]
	 Dùng các ký hiệu $\forall$, $\exists$ trước các mệnh đề chứa biến để được mệnh đề đúng
	\begin{multicols}{3}
		\begin{enumerate}[a.]
			\item $x+2>3$. 
			\item $a+3=3+a$. 
			\item  $15$ là bội của $x$.
			\item $\left(x-2\right)^2>-1$. 
			 \item $x+y>1$.
			\item $(a-b)(a+b)=a^2-b^2$. 
			\item $\left(a-b\right)^2=a^2-b^2$. 
			 \item $x^2>0$.
			 \item $\left(x+y\right)^2=x^2+2xy+y^2$.
			\item $\left(x-2\right)^2=1$. 
			 \item $x^2-5x+6=0$.
			 \item $(x+y)z=xz+yz$.
		\end{enumerate}
	\end{multicols}
	 \loigiai{
		\begin{multicols}{2}
			\begin{enumerate}[a.]
				\item $\exists x\in  \mathbb{R}:x+2>3$.
				 \item $\forall a\in \mathbb{R}:a+3=3+a$.
				\item  $\exists x\in \mathbb{R}:15$ là bội của $x$.
				\item  $\forall x\in \mathbb{R}:\left(x-2\right)^2>-1$.
				\item  $\exists x,y\in \mathbb{R}:x+y>1$.
				\item $\forall  a,b\in \mathbb{R}:(a-b)(a+b)=a^2-b^2$.
				\item $\exists  a,b\in \mathbb{R}:\left(a-b\right)^2=a^2-b^2$.
				\item $\exists x\in \mathbb{R}:x^2>0$. 
				\item  $\forall x,y\in \mathbb{R}:\left(x+y\right)^2=x^2+2xy+y^2$.
				\item $\exists  x\in \mathbb{R}:\left(x-2\right)^2=1$.
				\item $\exists x\in \mathbb{R}:x^2-5x+6=0$. 
				\item  $\forall x,y,z\in \mathbb{R}:(x+y)z=xz+yz$.
			\end{enumerate}
		\end{multicols}
	 }
\end{bt}

% nguồn: MaHoaNH_TuLuan_Lop10
\begin{bt}%[0D1B1-5]%[ma:3a32ef53]
	Nêu mệnh đề phủ định của các mệnh đề sau 
	\begin{enumEX}{2}
		\item $\forall x\in  \mathbb{R}:x^2>0$.
		\item $\exists x\in  \mathbb{R}:x>x^2$.
		\item $\exists x\in \mathbb{Q}:4x^2-1=0$. 
		\item  $\forall x\in \mathbb{R}:x^2-x+7>0$.
		\item $\forall x\in  \mathbb{R}:x^2-x-2<0$.
		\item $\exists  x\in \mathbb{R}:x^2=3$.
		\item $\forall  n\in \mathbb{N}$, $n^2+1$ không chia hết cho $3$.
		\item  $\forall n\in \mathbb{N}$, $n^2+2n+5$ là số nguyên tố.
		\item $\forall n\in \mathbb{N}$, $n^2+n$ chia hết cho $2$. 
		\item $\forall n\in \mathbb{N}$, $n^2-1$ là số  lẻ.
	\end{enumEX}
	 \loigiai{
		\begin{enumEX}{1}
			\item  $\forall x\in \mathbb{R}:x^2>0$.\\
			Mệnh đề phủ định \lq\lq$\exists x\in \mathbb{R}:x^2\le 0$\rq\rq. 
			\item $\exists  x\in \mathbb{R}:x>x^2$.\\
			 Mệnh đề phủ định \lq\lq $\forall x\in \mathbb{R}:x\le x^2$\rq\rq.
			 \item $\exists x\in \mathbb{Q}:4x^2-1=0$.\\
			Mệnh đề phủ định \lq\lq $\forall x\in \mathbb{Q}:4x^2-1\ne  0$\rq\rq.
			 \item $\forall x\in \mathbb{R}:x^2-x+7>0$.\\
			Mệnh đề phủ định \lq\lq $\exists x\in \mathbb{R}:x^2-x+7\le 0$\rq\rq. 
			\item $\forall x\in \mathbb{R}:x^2-x-2<0$.\\ 
			Mệnh đề phủ định \lq\lq $\exists x\in  \mathbb{R}:x^2-x-2\ge 0$\rq\rq.
			\item $\exists  x\in \mathbb{R}:x^2=3$.\\
			Mệnh  đề phủ định \lq\lq $\forall x\in \mathbb{R}:x^2\ne 3$\rq\rq.
			\item $\forall n\in  \mathbb{N}$, $n^2+1$ không chia hết cho $3$.\\
			Mệnh đề phủ định \lq\lq $\exists n\in \mathbb{N}$,  $n^2+1$ chia hết cho $3$\rq\rq.
			\item $\forall n\in \mathbb{N}$, $n^2+2n+5$ là  số nguyên tố.\\
			Mệnh đề phủ định \lq\lq $\exists  n\in \mathbb{N}$, $n^2+2n+5$ không là số nguyên tố\rq\rq.
			\item $\forall n\in \mathbb{N}$, $n^2+n$ chia hết cho  $2$.\\
			Mệnh đề phủ định  \lq\lq $\exists n\in \mathbb{N}$, $n^2+n$ không chia hết cho $2$\rq\rq.
			\item $\forall n\in \mathbb{N}$,  $n^2-1$ là số lẻ.\\
			Mệnh đề phủ định \lq\lq $\exists n\in \mathbb{N}$, $n^2-1$ là  số chẵn\rq\rq.
		\end{enumEX}
	 }
\end{bt}

% nguồn: MaHoaNH_TuLuan_Lop10
\begin{bt}%[0D1B1-5]%[ma:ddcc8dd0]
	Các mệnh đề sau đây đúng hay  sai? Giải thích? Viết mệnh đề phủ định của chúng?
	\begin{listEX}[2]
		 \item $\exists n\in \mathbb{Q}$, $n^2=2$.
		 \item $\forall x\in \mathbb{R}$, $x>x^2$.
		\item $\exists x\in \mathbb{R}$,  $x>x^2$.
		\item  $\forall n\in \mathbb{N}$, $n^2\ge n$.
		 \item $\exists n\in \mathbb{N}$, $n^2\ge n$.
		\item  $\forall x\in \mathbb{R}$, $x^2-x+1>0$.
		\item $\exists x\in \mathbb{R}$, $x^2-x+1>0$. 
		\item $\forall n\in \mathbb{N}$, $n^2+1$ không chia  hết cho $3$.
		\item  $\exists n\in \mathbb{N}$, $n^2+1$ không chia hết cho $3$.
		\item $\exists n\in \mathbb{N}$,  $n^2+1$ chia hết cho $4$.
	\end{listEX}
	\loigiai{ 
		\begin{listEX}
			\item $\exists n\in  \mathbb{Q}$, $n^2=2$.\\
			Mệnh đề sai vì phương trình $n^2=2$ không có nghiệm trên tập hợp $\mathbb{Q}$. Mệnh  đề phủ định là \lq\lq $\forall n\in \mathbb{Q}$, $n^2\ne 2$\rq\rq.
			\item $\forall  x\in \mathbb{R}$, $x>x^2$.\\
			Mệnh đề sai vì $x=0$ không thỏa mãn. Mệnh đề phủ định là \lq\lq $\exists  x\in \mathbb{R}$, $x\le x^2$\rq\rq.
			\item $\exists x\in \mathbb{R}$, $x>x^2$.\\ 
			Mệnh đề đúng vì $x=\dfrac{1}{2}$ thỏa mãn.  Mệnh đề phủ định là \lq\lq $\forall x\in \mathbb{R}$, $x\le x^2$\rq\rq.
			 \item $\forall n\in \mathbb{N}$, $n^2\ge n$.\\
			Mệnh đề vì đúng vì $n=0$ đúng và $n(n-1)\ge 0$, $\forall n\in \mathbb{N^*}$. Mệnh đề phủ định là \lq\lq $\exists n\in \mathbb{N}$, $n^2< n$\rq\rq. 
			\item $\exists n\in \mathbb{N}$, $n^2\ge n$.\\ 
			Mệnh đề đúng vì $n=0$ thỏa mãn. Mệnh đề phủ định là \lq\lq $\forall  n\in \mathbb{N}$, $n^2< n$\rq\rq.
			\item  $\forall x\in \mathbb{R}$, $x^2-x+1>0$.\\
			Mệnh đề đúng vì $x^2-x+1=\left( x-\dfrac{1}{2}\right)^2+\dfrac{3}{4}>0$, $\forall x\in \mathbb{R}$. Mệnh đề phủ định là  \lq\lq $\exists x\in \mathbb{R}$, $x^2-x+1\le 0$\rq\rq.
			\item $\exists x\in \mathbb{R}$, $x^2-x+1>0$.\\ 
			Mệnh đề đúng vì $x=0$ thỏa mãn. Mệnh đề phủ định là \lq\lq $\forall x\in  \mathbb{R}$, $x^2-x+1\le 0$\rq\rq.
			\item $\forall n\in  \mathbb{N}$, $n^2+1$ không chia hết cho $3$.\\
			Mệnh đề đúng. Ta  chứng minh như sau.\\
			Ta có một cách biểu diễn số tự nhiên $n$ là  $n=3k+r$ (với $k\in \mathbb{N}$ và $r\in \{ 0;1;2\}$).\\
			Nếu $r=0$, ta có  $n^2+1=9k^2+1$ không chia hết cho $3$.\\
			Nếu $r=1$, ta có $n^2+1=9k^2+6k+2$ không chia hết cho  $3$.\\
			Nếu $r=2$, ta có $n^2+1=9k^2+12k+4$  không chia hết cho $3$.\\
			Mệnh đề phủ định là \lq\lq  $\exists n\in \mathbb{N}$, $n^2+1$ chia hết cho $3$\rq\rq.
			\item $\exists n\in  \mathbb{N}$, $n^2+1$ không chia hết cho $3$.\\
			Mệnh đề đúng vì $n=0$ thỏa mãn. Mệnh đề phủ định là \lq\lq $\forall n\in  \mathbb{N}$, $n^2+1$ chia hết cho $3$\rq\rq.
			\item $\exists n\in \mathbb{N}$,  $n^2+1$ chia hết cho $4$.\\
			Mệnh đề sai.  Ta chứng minh như sau.\\
			Nếu $n$ là số chẵn, nghĩa là $n=2k$, $k\in \mathbb{N}$. Suy ra $n^2+1=4k^2+1$ không chia hết cho  $4$.\\
			Nếu $n$ là số lẻ, nghĩa là $n=2k+1$,  $k\in \mathbb{N}$. Suy ra $n^2+1=4k^2+4k+2$ không chia hết cho $4$.\\
			Vậy $\forall n\in \mathbb{N}$, $n^2+1$ không chia hết  cho $4$.\\
			Mệnh đề  phủ định là \lq\lq $\forall n\in \mathbb{N}$, $n^2+1$ không chia hết cho $4$\rq\rq.
		\end{listEX}
	} 
\end{bt}

% nguồn: MaHoaNH_TuLuan_Lop10
\begin{bt}%[0D1B1-5]%[ma:ea7f7a91]
	Trong các mệnh đề sau, mệnh đề nào đúng? Giải  thích? Phát biểu các mệnh đề đó thành lời?
	\begin{enumEX}{2}
		\item $\forall x\in \mathbb{R}$, $x^2>0$. 
		 \item $\exists x\in \mathbb{R}$, $x>x^2$.
		\item $\exists  x\in \mathbb{Q}$, $4x^2-1=0$.
		\item $\forall x\in  \mathbb{N}$, $n^2>n$.
		\item $\forall x\in \mathbb{R}$,  $x^2-x=1>0$.
		\item $\forall x\in \mathbb{R}$,  $x^2>9\Rightarrow x>3$.
		\item $\forall  x\in \mathbb{R}$, $x>3\Rightarrow x^2>9$.
		\item $\forall x\in \mathbb{R}$, $x^2<5\Rightarrow x<\sqrt{5}$. 
		\item $\exists x\in  \mathbb{R}$, $5x-3x^2\le 1$.
		\item $\exists x\in \mathbb{R}$, $x^2+2x+5$ là  hợp số.
		\item $\forall x\in \mathbb{N}$, $n^2+1$ không chia hết cho  $3$.
		 \item $\forall x\in \mathbb{N^*}$, $n(n+1)$ là số lẻ.
		\item $\forall  x\in \mathbb{N^*}$, $n(n+1)(n+2)$ chia hết cho $6$.
		 \item $\forall x\in \mathbb{N^*}$, $n^3+11n$ chia hết cho $6$.
	\end{enumEX}
	 \loigiai{
		\begin{enumEX}{1}
			\item  $\forall x\in \mathbb{R}$, $x^2>0$.\\
			Mệnh đề sai vì $x=0$ không thỏa mãn. Phát biểu \lq\lq Bình phương của một số  thực bất kì đều lớn hơn $0$\rq\rq.
			\item $\exists x\in \mathbb{R}$, $x>x^2$.\\ 
			Mệnh đề đúng vì $x=0{,}5$ thỏa mãn. Phát biểu \lq\lq Tồn tại một số thực lớn hơn bình phương của chính  nó\rq\rq.
			 \item $\exists x\in \mathbb{Q}$, $4x^2-1=0$.\\
			Mệnh đề đúng vì  $x=\dfrac{1}{2}$ thỏa mãn. Phát biểu \lq\lq Tồn tại một số hữu tỉ là nghiệm của phương trình $4x^2-1=0$\rq\rq.
			\item $\forall  n\in \mathbb{N}$, $n^2>n$.\\
			Mệnh đề  sai vì $n=0$ không thỏa mãn. Phát biểu \lq\lq Bình phương của một số tự nhiên bất kì đều lớn hơn chính nó\rq\rq.
			\item $\forall x\in \mathbb{R}$,  $x^2-x-1>0$.\\
			Mệnh đề sai vì $x=0$ không thỏa mãn. Phát biểu \lq\lq  Tất cả số thực đều thỏa mãn $x^2-x-1>0$\rq\rq.
			\item  $\forall x\in \mathbb{R}$, $x^2>9\Rightarrow x>3$.\\
			Mệnh đề sai vì với $x=-4$ thì $x^2>9$ đúng nhưng $x>3$ sai. Phát biểu \lq\lq Nếu bình phương một số thực lớn  hơn $9$ thì số đó lớn hơn $3$\rq\rq.
			\item  $\forall x\in \mathbb{R}$, $x>3\Rightarrow x^2>9$.\\
			Mệnh đề đúng. Phát biểu \lq\lq Nếu một số thực lớn hơn $3$  thì bình phương của nó lớn hơn $9$\rq\rq.
			\item $\forall x\in \mathbb{R}$, $x^2<5\Rightarrow  x<\sqrt{5}$.\\
			Ta có $x^2<5\Leftrightarrow (x-\sqrt{5})(x+\sqrt{5})<0\Leftrightarrow \heva{&x<\sqrt{5}\\&x>-\sqrt{5}}\text{ hoặc }\heva{&x>\sqrt{5}\\&x<-\sqrt{5}}\Leftrightarrow -\sqrt{5}<x<\sqrt{5}$.\\ 
			Mệnh đề đúng. Phát biểu \lq\lq Nếu bình phương một số thực  bé hơn $5$ thì nó bé hơn $\sqrt{5}$\rq\rq.
			\item $\exists x\in  \mathbb{R}$, $5x-3x^2\le 1$.\\
			Mệnh đề đúng vì $x=0$ thỏa mãn. Phát biểu \lq\lq  Tồn tại một số thực thỏa mãn $5x-3x^2\le 1$\rq\rq.
			\item $\exists  x\in \mathbb{R}$, $x^2+2x+5$ là hợp số.\\
			Mệnh đề đúng vì $x=1$ thỏa mãn. Phát biểu \lq\lq  Có một số thực thỏa mãn $x^2+2x+5$ là hợp số\rq\rq.
			\item $\forall n\in \mathbb{N}$, $n^2+1$  không chia hết cho $3$.\\
			Mệnh đề đúng. Ta chứng  minh như sau.\\
			Ta  có một cách biểu diễn số tự nhiên $n$ là $n=3k+r$ (với $k\in \mathbb{N}$ và $r\in \{ 0;1;2\}$).\\
			Nếu $r=0$, ta có $n^2+1=9k^2+1$ không chia hết cho $3$.\\ 
			Nếu $r=1$, ta có $n^2+1=9k^2+6k+2$ không chia hết  cho $3$.\\
			 Nếu $r=2$, ta có $n^2+1=9k^2+12k+4$ không chia hết cho $3$.\\
			Phát biểu \lq\lq Mọi số tự nhiên bình phương lên và cộng  thêm $1$ đều không chia hết cho $3$\rq\rq.
			\item  $\forall n\in \mathbb{N^*}$, $n(n+1)$ là số lẻ.\\
			Mệnh đề sai vì $n$ và $n+1$ là hai số tự nhiên liên tiếp nên $n(n+1)$ chia hết cho $2$. Vậy $n(n+1)$ là số chẵn. Phát  biểu \lq\lq Tích hai số tự nhiên liên tiếp khác $0$ đều là số lẻ\rq\rq.
			\item $\forall  n\in \mathbb{N^*}$, $n(n+1)(n+2)$ chia hết cho $6$.\\
			Mệnh đề đúng vì $n$, $n+1$, $n+2$ là ba số tự nhiên liên  tiếp nên luôn có một số chia hết cho $2$ và một số chia hết cho $3$. Vậy $n(n+1)(n+2)$ chia hết cho $6$. Phát biểu \lq\lq Tích ba số tự nhiên liên tiếp khác $0$ đều chia hết cho $6$\rq\rq.
			\item  $\forall n\in \mathbb{N^*}$, $n^3+11n$ chia hết cho $6$.\\
			Mệnh đề đúng  vì $n^3+11n=n^3-n+12n=(n-1)n(n+1)+12n$ chia hết cho $6$. Phát biểu \lq\lq Mọi số tự nhiên $n\ne 0$ đều thỏa $n^3+11n$ chia hết cho $6$\rq\rq.
		\end{enumEX}
	 }
\end{bt}
